\documentclass[11pt]{article}

\usepackage[T1]{fontenc}
\usepackage[utf8]{inputenc}
\usepackage[english]{babel}
\usepackage[margin=1in]{geometry}
\usepackage{lmodern}
\usepackage{microtype}
\usepackage{mathtools}
\usepackage{amssymb}
\usepackage{amsthm}
\usepackage{graphicx}
\usepackage{booktabs}
\usepackage{longtable}
\usepackage{array}
\usepackage{pdflscape}
\usepackage{caption}
\usepackage{enumitem}
\usepackage{fancyhdr}
\usepackage[numbers,sort&compress]{natbib}
% orcidlink: use package if available, else provide a fallback
\IfFileExists{orcidlink.sty}{\usepackage{orcidlink}}{%
  \newcommand{\orcidlink}[1]{\textsuperscript{\href{https://orcid.org/#1}{ORCID}}}%
}
\usepackage{titling}
\usepackage{doi}
\usepackage{xurl}
\usepackage{hyperref}

\hypersetup{
  pdftitle={Six Birds Foundations II: Admissibility, Meta-Theory, and the Exact-Six Program},
  pdfauthor={Ioannis Tsiokos},
  hidelinks
}
\urlstyle{same}
\captionsetup{
  font=small,
  labelfont=bf,
  labelsep=period,
  margin=1em
}
\captionsetup[table]{skip=6pt,position=top}
\captionsetup[figure]{skip=6pt}
\setlist[enumerate]{leftmargin=*, itemsep=2pt, topsep=4pt}
\setlist[itemize]{leftmargin=*, itemsep=2pt, topsep=4pt}
\setlength{\droptitle}{-3.2em}
\pretitle{\begin{center}\LARGE\bfseries}
\posttitle{\par\end{center}\vspace{0.5em}}
\preauthor{\begin{center}\normalsize}
\postauthor{\par\end{center}\vspace{-0.4em}}
\predate{\begin{center}\small}
\postdate{\par\end{center}\vspace{0.6em}}
\setlength{\emergencystretch}{2em}
\clubpenalty=10000
\widowpenalty=10000
\displaywidowpenalty=10000
\interfootnotelinepenalty=10000
\allowdisplaybreaks[2]
\setlength{\bibsep}{2pt plus 0.3ex}

\pagestyle{fancy}
\fancyhf{}
\fancyhead[L]{\small\itshape Six Birds Foundations II}
\fancyhead[R]{\small\thepage}
\renewcommand{\headrulewidth}{0.4pt}
\renewcommand{\footrulewidth}{0pt}

\fancypagestyle{firstpage}{\fancyhf{}
  \fancyfoot[C]{\scriptsize Automorph Inc., Wilmington, DE, USA \quad \texttt{ioannis@automorph.io} \quad \textcopyright\ 2026 \quad CC-BY 4.0 \quad Preprint --- not peer reviewed.}
  \renewcommand{\headrulewidth}{0pt}
  \renewcommand{\footrulewidth}{0pt}
}

\usepackage{titlesec}
\titlespacing*{\section}{0pt}{2.8ex plus 1.0ex minus 0.4ex}{1.6ex plus 0.4ex}
\titlespacing*{\subsection}{0pt}{2.2ex plus 0.8ex minus 0.3ex}{1.1ex plus 0.3ex}

\newcommand{\classid}[1]{\texttt{\small\hyphenchar\font=`\_ #1}}

\newtheorem{theorem}{Theorem}[section]
\newtheorem{proposition}[theorem]{Proposition}
\newtheorem{lemma}[theorem]{Lemma}
\newtheorem{corollary}[theorem]{Corollary}
\theoremstyle{definition}
\newtheorem{definition}[theorem]{Definition}
\newtheorem{remark}[theorem]{Remark}

\newcommand{\R}{\mathbb{R}}
\newcommand{\N}{\mathbb{N}}

\newcommand{\FATCD}{\ensuremath{\mathsf{FATCD}}}
\newcommand{\Raw}[1]{\ensuremath{\operatorname{Raw}\!\left(#1\right)}}
\newcommand{\Audit}[1]{\ensuremath{\operatorname{Audit}\!\left(#1\right)}}
\newcommand{\Closed}[1]{\ensuremath{\operatorname{Closed}\!\left(#1\right)}}
\newcommand{\Dec}[1]{\ensuremath{\operatorname{Dec}\!\left(#1\right)}}
\newcommand{\Residual}[1]{\ensuremath{\operatorname{Residual}\!\left(#1\right)}}
\newcommand{\Chan}[2]{\ensuremath{\operatorname{Chan}_{#1}\!\left(#2\right)}}
\newcommand{\Status}[2]{\ensuremath{\operatorname{Status}_{#1}\!\left(#2\right)}}
\newcommand{\ScopedExactSix}[1]{\ensuremath{\operatorname{ScopedExactSix}\!\left(#1\right)}}

\newcommand{\activeproj}{\texttt{active\_projection}}
\newcommand{\belowthr}{\texttt{below\_threshold}}
\newcommand{\collapsedbc}{\texttt{collapsed\_boundary\_case}}
\newcommand{\absentrec}{\texttt{absent\_with\_record}}
\newcommand{\outsidescope}{\texttt{inapplicable\_outside\_scope}}

\newcommand{\Pone}{\ensuremath{\mathrm{P1}}}
\newcommand{\Ptwo}{\ensuremath{\mathrm{P2}}}
\newcommand{\Pthree}{\ensuremath{\mathrm{P3}}}
\newcommand{\Pfour}{\ensuremath{\mathrm{P4}}}
\newcommand{\Pfive}{\ensuremath{\mathrm{P5}}}
\newcommand{\Psix}{\ensuremath{\mathrm{P6}}}
 
\begin{document}

\title{Six Birds Foundations II:\\[4pt]
Admissibility, Meta-Theory, and the Exact-Six Program}
\author{Ioannis Tsiokos\,\orcidlink{0009-0009-7659-5964}}
\date{\today}

\maketitle
\thispagestyle{firstpage}

\begin{abstract}
\noindent
Six Birds Foundations I~\citep{Tsiokos2026SixBirdsFoundations}
introduced a stratified closure calculus with six primitive closure
roles. The present paper extends that calculus
into a theorem-level framework in three parts: a primitive role
architecture fixing the meanings, levels, local boundaries, and
selected lifts of the six typed closure roles \(P_1,\ldots,P_6\); an
admissibility and meta-theory regime in which honest bookkeeping,
typed non-collapse, level profiles, forgetting with selected lifts,
activation thresholds, instrument-relative visibility, and
empirical-bridge admissibility are stated as theorem schemas; and a
scoped exact-six theorem chain over the domain \(\FATCD{}\) of finite
audited typed closure descriptions. Within this domain we prove six
role-lower-bound witnesses, an upper-bound classification with
corrected role-projection alignments (closing the probability,
causality, and agency threats in the covered scope), a
decomposition/exhaustion theorem, and an integrated stress
certification. The terminal theorem asserts that every
\(D \in \FATCD{}\) admits a well-formed decomposition/exhaustion
record with exactly six role channels, each in one of five recorded
statuses, with residual features classified under an explicit scheme
and, under an explicit recognition hypothesis (the three principal threat
classes --- probability, causality, and agency --- closed unconditionally), no
seventh irreducible role required in the covered scope. The
claim concerns role channels rather than active projections, and the
decomposition is existential, not unique. We do not claim unrestricted
exact-six, all-emergence coverage, a full six-primitives algebra,
global primitive independence, empirical realization, internal
instrument-family completeness, or recognition-completeness of the six roles.
\end{abstract}
 \section{Introduction}\label{sec:introduction}



Six Birds Foundations I~\citep{Tsiokos2026SixBirdsFoundations}
introduced a stratified closure-calculus
perspective~\citep{Anderson1972} on finite theory packages,
\(\mathcal{T} = (Z, f, \Sigma_f, E, \mathcal{A})\), in which theories
arise as fixed points of idempotent
operators~\citep{DaveyPriestley2002} and six primitive closure roles
\Pone{}, \ldots, \Psix{} emerge from composability combined with
limited-interface observation. The present paper extends
that calculus into a theorem-level framework. It fixes the primitive
role architecture, specifies the admissibility and meta-theory regime
under which role talk is licensed, and proves a scoped exact-six
terminal result over a precisely delimited domain. The exact-six
theorem is the consequence of that stack, not its whole content.

\subsection{The three parts}\label{subsec:three-pillars}

Each of the three parts below supplies structure the other two do not.
The \emph{primitive role architecture}
(Section~\ref{sec:primitive-role-architecture}) fixes the six role
meanings of \Pone{}--\Psix{} at their corrected alignments, records
the common behavioral, verification, and structural-categorical level
pattern, enumerates the local boundary distinctions under scoped
bookkeeping, and states the selected-lift atlas vocabulary. The
\emph{admissibility and meta-theory regime}
(Section~\ref{sec:admissibility-meta-theory}) elevates honest
bookkeeping from a per-description condition to a meta-theoretic
discipline, replaces flat primitive independence by typed
non-collapse, and records level, forgetting, threshold, visibility,
and empirical-bridge theorem schemas. The \emph{exact-six theorem
chain}
(Sections~\ref{sec:fatcd-domain}--\ref{sec:integrated-stress})
quantifies the exact-six question precisely over the domain \FATCD{}
of finite audited typed closure descriptions and proves the terminal
result via lower-bound witnesses, an upper-bound classification with
corrected role-projection alignments, a decomposition/exhaustion
theorem, and an integrated stress certification.

\subsection{The terminal theorem, informally}\label{subsec:terminal-informal}

A finite audited typed closure description \(D\) is a finitely
presented collection of typed raw records and annotations subject to
finiteness, closure, and honest-bookkeeping conditions. The role
vocabulary \Pone{}--\Psix{} enters \(D\) not as a constraint on
admissibility but as a family of derived projections \(\pi_i(D)\),
each available when threshold, witness, level, and audit data are
attached to raw features. The terminal theorem
(Theorem~\ref{thm:scoped-exact-six}) states that for every
\(D \in \FATCD{}\) there exists a well-formed
decomposition/exhaustion record \(\Dec{D}\) carrying exactly the six
typed closure-role channels \Pone{}, \ldots, \Psix{}, each in exactly
one of five statuses, with all residual features covered by an
explicit classification and, under an explicit recognition hypothesis, no
seventh irreducible role required in the covered scope. The claim is about role channels, not active
projections, and the decomposition is existential, not unique.

\subsection{What the paper does and does not claim}\label{subsec:intro-nonclaims}

The terminal theorem is scoped exact-six, not unrestricted exact-six.
It makes none of the following stronger claims:
\begin{itemize}[topsep=2pt,itemsep=1pt]
  \item unrestricted exact-six, or any claim about all emergence
    phenomena;
  \item that every \(D \in \FATCD{}\) activates all six roles or
    carries six active nontrivial witnesses;
  \item that the decomposition/exhaustion record \(\Dec{D}\) is unique
    or canonical;
  \item a full six-primitives algebra or global primitive
    independence;
  \item empirical realization;
  \item internal instrument-family completeness;
  \item any statement about broader domains without further work.
\end{itemize}
These are scope boundaries, collected in
Section~\ref{sec:limitations-nonclaims} and enforced locally by the
theorem-chain results that would otherwise overread them.

\subsection{Organization}\label{subsec:organization}

After this introduction, Section~\ref{sec:primitive-role-architecture}
fixes the primitive role meanings, level trichotomy, boundary
distinctions, and selected-lift atlas. Section~\ref{sec:admissibility-meta-theory}
states the admissibility and meta-theory regime. Section~\ref{sec:fatcd-domain}
defines \FATCD{} and formulates the exact-six question over it.
Section~\ref{sec:lower-bounds} proves the six-role lower bounds.
Section~\ref{sec:upper-bound-classification} proves the repaired
upper-bound classification theorem, closing the probability,
causality, and agency threats within the covered scope.
Section~\ref{sec:decomposition-exhaustion} proves decomposition
existence, six-channel exactness, and residual exhaustion, and records
the channel-versus-activation discipline and non-uniqueness.
Section~\ref{sec:integrated-stress} records the anti-overclaim stress
families and the integrated stress certification.
Section~\ref{sec:scoped-exact-six} states and proves the terminal
theorem. Sections~\ref{sec:limitations-nonclaims}
and~\ref{sec:future-work} record scope boundaries and future
directions. Appendix~\ref{app:provenance} records development provenance and
evidence hygiene, and Appendix~\ref{app:mechanization} documents the
Lean mechanization.
 \section{Primitive role architecture}\label{sec:primitive-role-architecture}



This section fixes the meanings of the six typed closure roles
\Pone{}, \ldots, \Psix{}, the behavioral, verification, and
structural-categorical level pattern they share, the local boundary
distinctions between them, and the selected-lift atlas vocabulary.
These meanings are fixed for the remainder of the paper and govern
every subsequent definition, lower bound, and upper-bound
classification.

\subsection{Role meanings}\label{subsec:role-meanings}

\begin{definition}[Primitive closure roles]\label{def:primitive-roles}
The six primitive closure roles are
\begin{itemize}[topsep=2pt,itemsep=1pt]
  \item \Pone{} --- update-vs-quotient compatibility / descent
    obstruction~\citep{GrothendieckSGA1,JanelidzeTholen1994};
  \item \Ptwo{} --- macro-specification
    representability~\citep{MacLane1998};
  \item \Pthree{} --- enriched route/coherence
    mismatch~\citep{BaaderNipkow1998,Nakahara2003}; \Pthree{} is not a
    directionality certificate;
  \item \Pfour{} --- stage/order/transition/refinement
    structure~\citep{Milner1989,BackVonWright1998};
  \item \Pfive{} --- quotient packaging /
    quotienting~\citep{KemenySnell1960};
  \item \Psix{} --- audit/bookkeeping/provenance/replay/checkability~\citep{GreenKarvounarakisTannen2007,Merkle1987}.
\end{itemize}
\end{definition}

\begin{remark}[Corrected alignments]\label{rem:repaired-alignments}
Three alignments that Definition~\ref{def:primitive-roles} rules out
are used throughout later classification: generic transition data is
not \Pthree{}; generic path or derivation data is not \Pfive{}; and
generic semantic or interpretation data is not \Pone{}. These
constraints determine which feature clusters can later be interpreted
as active projections of each role.
\end{remark}

\begin{remark}[Drive face of \Psix{}]\label{rem:p6-drive}
Arguments about affinities, detailed balance, or honest directionality
use the narrower \(\mathrm{P6}_{\mathrm{drive}}\) specialization of
\Psix{}: its affinity-carrying, detailed-balance-breaking
face~\citep{Seifert2012}. \(\mathrm{P6}_{\mathrm{drive}}\) is a face
of \Psix{}, not a replacement for it. A \Pthree{}
route/coherence-mismatch witness does not establish directionality
without a \Psix{}-side audit.
\end{remark}

\subsection{Behavioral, verification, and structural-categorical levels}\label{subsec:levels}

Each role admits a common three-level pattern.

\begin{definition}[Level trichotomy]\label{def:level-trichotomy}
For each role \Pone{}, \ldots, \Psix{}, we distinguish
\begin{itemize}[topsep=2pt,itemsep=1pt]
  \item the \emph{behavioral level}: status, occurrence, or observed
    behavior of the role in a description;
  \item the \emph{verification level}: proof, witness, certificate, or
    checkable record supporting a behavioral
    claim~\citep{MartinLof1984};
  \item the \emph{structural-categorical level}: reusable structure
    supporting the primitive role.
\end{itemize}
Forgetting proceeds from structural-categorical to verification to
behavioral and is lossy. Lifting proceeds from behavioral to
verification to structural-categorical only after adding explicitly
recorded data, and selected lifts are not unique.
\end{definition}

\begin{remark}[Level content per primitive]\label{rem:per-primitive-levels}
The level trichotomy instantiates as follows.
\begin{itemize}[topsep=2pt,itemsep=1pt]
  \item \Pone{}: descent success or descent obstruction status; descent
    proof or obstruction witness; quotient data equipped with a descended
    map, or an explicit no-descended-map obstruction record.
  \item \Ptwo{}: representability or non-representability status;
    representability or non-representability certificate;
    macro-specification layer together with a realization relation.
  \item \Pthree{}: witness behavior of route disagreement; generated route
    terms and non-reducibility certificates; categorical route grammar or a
    reduced categorical universe.
  \item \Pfour{}: stage status, transition occurrence, or refinement
    behavior; order/preorder, transition, or refinement certificates;
    stage/refinement architecture.
  \item \Pfive{}: packaging or identification behavior; equivalence and
    quotient-validity certificate; quotient object with a packaging map
    and universal-property or tower/coarsening data.
  \item \Psix{}: audit-relevant event occurrence or record--event
    association; checkable audit record; audit schema with provenance,
    replay, and check relations.
\end{itemize}
The meta-theoretic properties of this trichotomy --- lossy forgetting,
selected non-unique lifts, and the absence of lower-to-higher determinacy
without added data --- are formalized in
Section~\ref{sec:admissibility-meta-theory}.
\end{remark}

\subsection{Local boundary distinctions}\label{subsec:boundary-matrix}

With role meanings and level pattern fixed, we turn to how the six
roles remain distinct. Table~\ref{tab:boundary-matrix} records the
pairwise local rules. These are local distinctions under scoped
bookkeeping, not a global primitive-independence theorem: two roles
may interact legitimately in a single description, but neither
collapses into the other. The table is reproduced in landscape
orientation on the page immediately preceding the references.

\begin{proposition}[Boundary distinctions under scoped bookkeeping]\label{prop:boundary-distinctions}
For every pair of roles listed in Table~\ref{tab:boundary-matrix}, the
corresponding ``forbidden collapse'' is ruled out as a locally admissible
identification of those two roles under the repaired alignments of
Remark~\ref{rem:repaired-alignments} together with the honest bookkeeping
required by a \FATCD{}.\footnote{Tracked in the formalization harness as a conditional schema (\texttt{SixBirds.boundary\_distinctions}); see Appendix~\ref{app:mechanization}.}
\end{proposition}

\begin{proof}
We argue row by row. Each row of Table~\ref{tab:boundary-matrix} names an
ordered pair \((P_i,P_j)\), a permitted interaction, and a forbidden
collapse. The forbidden collapse always has the same logical form: it would
license treating witness or structural data carried under the meaning of one
role as if it were, by itself, witness or structural data for the other role
--- that is, it would identify the two roles' admissible data types. We
establish the distinction by a general schema, then prove that the schema
provably applies to \emph{every} row by exhibiting the common form of its
forbidden-collapse column, and finally discharge all fifteen rows of
Table~\ref{tab:boundary-matrix} by a compact enumeration that maps each row's
forbidden collapse to the schema, with the role-specific data types read off
Definition~\ref{def:primitive-roles} and the per-primitive level content of
Remark~\ref{rem:per-primitive-levels}. Three rows are then worked in full as
illustrations; they add nothing the enumeration leaves open.

\emph{Schema.} Fix a row \((P_i,P_j)\). By
Definition~\ref{def:primitive-roles} the two roles have fixed, distinct
meanings, and by Remark~\ref{rem:per-primitive-levels} each meaning prescribes
a distinct kind of behavioral status, verification witness, and structural
object. A forbidden collapse would assert that data admissible as a
\(P_j\)-witness (or \(P_j\)-structure) is, with no further data,
admissible as a \(P_i\)-witness (or \(P_i\)-structure), or conversely.
Suppose such an identification were locally admissible inside a \FATCD{}.
Honest bookkeeping requires every active projection to carry a role-aligned
feature cluster together with the witness, threshold, level, and audit data
appropriate to its role, and requires the channel status of each role to be
recorded on its own role-aligned content rather than borrowed from another
role. Under the collapse, the bookkeeping for one channel would be discharged
by content whose recorded role meaning is the other channel's. The two role
meanings being distinct and non-interchangeable, that recorded content does not
meet the activation conditions of the channel it would be credited to: the
required witness type is simply not present. The honest record must then carry
that channel as \belowthr{} or \collapsedbc{}, not as \activeproj{}, which
contradicts the supposed identification. Hence no such identification is
locally admissible. The permitted interaction of the same row is untouched,
because it composes \(P_i\)-content with \(P_j\)-content while keeping
each role's bookkeeping on its own content; composition is not identification.

\emph{The schema applies to every row.} It remains to verify that the
hypothesis of the schema --- that the forbidden collapse credits one channel on
content whose recorded role meaning is the other channel's --- holds for each of
the fifteen rows, and not merely for chosen representatives. We read this off
the structure of the forbidden-collapse column. Every entry in that column of
Table~\ref{tab:boundary-matrix} has one of two shapes, both instances of the
schema. \emph{(Carried form.)} The entry has the form ``\(X\) alone is not
\(P_i\)'' or ``\(X\) is not \(P_i\)'', where \(X\) is content named by the
distinction-rule column as the meaning of the partner role \(P_j\) (packaging,
stage index, quotient object, audit record, and so on). Such an entry is exactly
the assertion the schema forbids: it would credit the \(P_i\) channel on the
strength of \(P_j\)-content. \emph{(Non-implication form.)} The entry has the
form ``\(P_j\)-status \(S\) is not automatically \(P_i\)-status \(S'\)'' (a
representability failure read as a descent obstruction, a route
undefinedness read as a route mismatch, a non-representability read as a route
mismatch). Such an entry is the symmetric case named in the schema: it would
credit the \(P_i\) channel's status on the strength of a status belonging to
\(P_j\)'s meaning. In both shapes the credited channel \(P_i\) and the
content/status \(X\) (or \(S\)) carried under \(P_j\)'s meaning are fixed by the
two columns of the row, and by
Definition~\ref{def:primitive-roles} together with
Remark~\ref{rem:per-primitive-levels} the witness type required to activate
\(P_i\) is distinct from, and not supplied by, the \(P_j\)-content the collapse
offers. The schema therefore forces the honest \(P_i\)-status to \belowthr{} or
\collapsedbc{}, contradicting the identification. This holds for every row,
because every forbidden-collapse entry falls under one of the two shapes just
described; the only remaining task is to record, per row, which channel is
credited and which witness type is thereby missing.

\emph{Enumeration of all fifteen rows.} We list each row as
\(P_i/P_j\), name the credited channel, the borrowed content, and the
distinct witness type the credited channel requires but does not receive, so
that the schema applies with no residue. Throughout, ``the honest status is
\belowthr{}'' marks the case in which the credited channel's own witness type is
simply not present in the borrowed content, while ``the honest status is
\collapsedbc{}'' marks the case in which a candidate cluster of the partner
role's data is offered as the credited channel's witness and fails to
instantiate that witness schema.
\begin{enumerate}[topsep=2pt,itemsep=2pt,label=\textup{(\arabic*)}]
  \item \(\Pone{}/\Pfive{}\), ``packaging alone is not \Pone{}'': credits
    \Pone{} on \Pfive{} packaging content; the required descended-map witness
    or no-descended-map obstruction record is absent, so the honest \Pone{}
    status is \belowthr{}.
  \item \(\Pone{}/\Pfour{}\), ``stage index alone is not \Pone{}'': credits
    \Pone{} on \Pfour{} stage-index content; no descent or obstruction witness
    is thereby present, so the honest \Pone{} status is \belowthr{}.
  \item \(\Pone{}/\Ptwo{}\), ``representability failure is not automatically
    descent obstruction'': credits \Pone{} on a \Ptwo{} non-representability
    status; a non-representability certificate is not a descended-map
    obstruction record and does not instantiate the \Pone{} obstruction
    witness, so the honest \Pone{} status is \collapsedbc{}.
  \item \(\Pone{}/\Pthree{}\), ``route undefinedness is not route mismatch'':
    credits \Pthree{} on a \Pone{} route-blocking status; a non-reducibility
    certificate between \emph{already-defined} routes is required, and an
    undefined route supplies no such certificate, so the candidate cluster does
    not instantiate the \Pthree{} witness schema and the honest \Pthree{} status
    is \collapsedbc{}.
  \item \(\Pone{}/\Psix{}\), ``audit record is not descent obstruction'':
    credits \Pone{} on a \Psix{} audit record; an audit record is not a
    descended-map or obstruction witness, so the honest \Pone{} status is
    \belowthr{}.
  \item \(\Ptwo{}/\Pfive{}\), ``nontrivial \Ptwo{} is not mere packaging'':
    credits \Ptwo{} on \Pfive{} packaging content in the nontrivial case; the
    realization relation beyond quotient equality that makes the \Ptwo{} content
    nontrivial is absent from the packaging map, so the honest \Ptwo{} status is
    \belowthr{}. (In the trivial case the row permits the coincidence and
    forbids nothing.)
  \item \(\Ptwo{}/\Pfour{}\), ``stage variation alone is not \Ptwo{}'':
    credits \Ptwo{} on \Pfour{} stage content; a representability certificate
    with realization relation is required, and stage variation supplies none, so
    the honest \Ptwo{} status is \belowthr{}.
  \item \(\Ptwo{}/\Pthree{}\), ``non-representability is not route mismatch'':
    credits \Pthree{} on a \Ptwo{} non-representability status; a
    non-representability record is not a non-reducibility certificate between
    defined routes and does not instantiate the \Pthree{} witness schema, so the
    honest \Pthree{} status is \collapsedbc{}.
  \item \(\Ptwo{}/\Psix{}\), ``audit trail is not representability'': credits
    \Ptwo{} on a \Psix{} audit trail; an audit record is not a representability
    certificate, so the honest \Ptwo{} status is \belowthr{}.
  \item \(\Pthree{}/\Pfour{}\), ``stage-indexed route mismatch is not \Pfour{}
    itself'': credits \Pfour{} on a \Pthree{} indexed route-disagreement
    witness; a non-reducibility certificate carries no order, transition, or
    refinement structure and does not instantiate the \Pfour{} witness schema,
    so the honest \Pfour{} status is \collapsedbc{}.
  \item \(\Pthree{}/\Pfive{}\), ``quotient object is not route mismatch'':
    credits \Pthree{} on a \Pfive{} quotient object; a quotient object is not a
    non-reducibility certificate between defined routes and does not instantiate
    the \Pthree{} witness schema, so the honest \Pthree{} status is
    \collapsedbc{}.
  \item \(\Pthree{}/\Psix{}\), ``audit record is not route mismatch'': credits
    \Pthree{} on a \Psix{} audit record; an audit record is not a
    non-reducibility certificate, so the honest \Pthree{} status is
    \belowthr{}.
  \item \(\Pfour{}/\Pfive{}\), ``packaging is not stage structure'': credits
    \Pfour{} on \Pfive{} packaging content; a packaging map is not an
    order/transition/refinement certificate and does not instantiate the
    \Pfour{} witness schema, so the honest \Pfour{} status is \collapsedbc{}.
  \item \(\Pfour{}/\Psix{}\), ``audit record is not staging'': credits \Pfour{}
    on a \Psix{} audit record of staged histories; an audit record is not an
    order/transition/refinement certificate, so the honest \Pfour{} status is
    \belowthr{}.
  \item \(\Pfive{}/\Psix{}\), ``audit record is not packaging map'': credits
    \Pfive{} on a \Psix{} audit record of packaging events; an audit record is
    not a quotient object with a packaging map, so the honest \Pfive{} status is
    \belowthr{}.
\end{enumerate}
Clauses (1)--(15) exhaust the rows of Table~\ref{tab:boundary-matrix}: each
names the credited channel, the borrowed partner-role content, and the
distinct witness type the credited channel requires and does not receive, and
in each the schema forces the honest status to \belowthr{} or \collapsedbc{}.
Nothing is left to follow ``mutatis mutandis''.

We now work three of these rows in full; they illustrate the carried and
non-implication shapes and add no obligation beyond clauses (1)--(15).

\emph{Row \(\Pone{}/\Pfive{}\) (packaging alone is not \(\Pone{}\)), clause
(1).} By Definition~\ref{def:primitive-roles}, \Pfive{} is quotient packaging
and \Pone{} is update-vs-quotient compatibility, i.e.\ descent. By
Remark~\ref{rem:per-primitive-levels} a \Pfive{}-witness is a quotient object
with a packaging map (and universal-property or tower data), whereas a
\Pone{}-witness is quotient data \emph{equipped with a descended map}, or else
an explicit no-descended-map obstruction record. Possession of the packaging
map is precisely the data that says \emph{what} the quotient is; it says nothing
about \emph{whether} a selected update descends through it. The collapse
``packaging alone is \Pone{}'' would credit the \Pone{} channel on the strength
of the packaging map alone, but no descended-map witness and no obstruction
record is thereby present, so by the schema the honest \Pone{} status is
\belowthr{}, not \activeproj{}: the descent question has not been answered. The
permitted interaction --- \Pone{} using \Pfive{} when checking descent or
obstruction --- is untouched, since there \Pone{} consumes the quotient
packaging as the object \emph{over which} a descended map or obstruction is then
exhibited, and the descent witness remains the \Pone{} channel's own content.

\emph{Row \(\Pthree{}/\Pfour{}\) (stage-indexed route mismatch is not
\(\Pfour{}\) itself), clause (10).} By Definition~\ref{def:primitive-roles},
\Pfour{} is stage/order/transition/refinement structure and \Pthree{} is
enriched route/coherence mismatch; by Remark~\ref{rem:repaired-alignments}
generic transition data is \emph{not} \Pthree{} and, symmetrically, a
route-mismatch witness is not \Pfour{}. A \Pfour{}-witness is an
order/transition/refinement certificate or a stage architecture; a
\Pthree{}-witness is a non-reducibility certificate exhibiting disagreement
between already-defined routes. When route mismatch is indexed by stages, the
stage index is \Pfour{}-content and the route disagreement at each index is
\Pthree{}-content; these are recorded as data of different roles. The collapse
``stage-indexed route mismatch is \Pfour{}'' would credit the \Pfour{} channel
on the strength of the indexed route-disagreement witness, but a
non-reducibility certificate is not an order/transition/refinement certificate
--- it carries no stage, order, or refinement structure of its own --- so by the
schema the \Pfour{} status it would license is \collapsedbc{}: the candidate
cluster does not instantiate the required \Pfour{} witness schema. The permitted
interaction --- route systems varying over valid stages --- is untouched, since
there the \Pfour{} stage structure is supplied independently and the \Pthree{}
mismatch is merely indexed by it, each channel keeping its own witness.

\emph{Row \(\Ptwo{}/\Pfive{}\) (nontrivial \(\Ptwo{}\) is not mere packaging),
clause (6).} By Definition~\ref{def:primitive-roles}, \Ptwo{} is
macro-specification representability and \Pfive{} is quotient packaging. As the
row itself records, when a specification is nothing more than a quotient class
the \Ptwo{} content \emph{may} legitimately coincide with \Pfive{} packaging;
the forbidden collapse is scoped to the \emph{nontrivial} case. By
Remark~\ref{rem:per-primitive-levels} a \Ptwo{}-witness is a representability
certificate for a macro-specification layer together with a realization
relation, while a \Pfive{}-witness is a quotient object with a packaging map.
Nontriviality of the \Ptwo{} content means exactly that its representability
certificate carries a realization relation not reducible to quotient equality
--- data beyond the packaging map. The collapse ``nontrivial \Ptwo{} is mere
packaging'' would credit the \Ptwo{} channel on the strength of the packaging
map alone; but the realization relation that makes the content nontrivial is by
hypothesis absent from the packaging data, so by the schema the honest \Ptwo{}
status is \belowthr{}: the representability witness required for activation is
not present. This argument is confined to the nontrivial case; in the trivial
case the row permits the coincidence and there is nothing to forbid. The
permitted interaction --- nontrivial \Ptwo{} carried over \Pfive{} quotient
structure --- is untouched, since there the realization relation is supplied as
the \Ptwo{} channel's own content layered over the \Pfive{} packaging.

The enumeration (1)--(15) covers every row of
Table~\ref{tab:boundary-matrix}, and the three worked rows confirm both shapes
in detail. The argument is local to each listed pair under the covered scope and
the honest bookkeeping a \FATCD{} requires; it establishes pairwise
non-collapse, not global independence of the six roles (cf.\
Remark~\ref{rem:not-global-independence}).
\end{proof}

\begin{remark}[Not a global independence theorem]\label{rem:not-global-independence}
The boundary matrix asserts local distinctions, not global
independence. Primitives interact legitimately in a single
description: \Pone{} descent is audited by \Psix{}; \Pthree{} route
terms may be indexed by \Pfour{} stages; nontrivial \Ptwo{}
representability is typically carried over \Pfive{} quotient structure.
No claim is made that the six roles are independent generators of a
global algebra. The typed-non-collapse counterpart of this observation
appears in Section~\ref{sec:admissibility-meta-theory}.
\end{remark}

\subsection{Selected-lift atlas and global nonclaims}\label{subsec:lift-atlas}

The remaining architectural data are the selected lifts between levels
and the nonclaims that accompany them. Upgrading from behavioral or
verification content to the structural-categorical level requires
explicitly recorded added data. The resulting lifts are selected, not
unique, and the atlas carries structural content that is lost under
forgetting.

\begin{definition}[Selected-lift atlas]\label{def:selected-lift-atlas}
The \emph{selected-lift atlas} for the six roles records, for each
\Pone{}, \ldots, \Psix{}:
\begin{itemize}[topsep=2pt,itemsep=1pt]
  \item \Pone{}: proof or witness together with quotient structure, carrying
    either a descended map or an explicit no-map obstruction;
  \item \Ptwo{}: representability or non-representability certificates
    together with the macro-specification layer and realization relation;
  \item \Pthree{}: the selected lift data recorded in the promoted atlas,
    carrying generated route terms, non-reducibility certificates, and
    categorical route grammar or a reduced categorical universe;
  \item \Pfour{}: validity certificates together with structural
    stage/refinement architecture;
  \item \Pfive{}: quotient-validity certificates together with the
    structural quotient data, a packaging map, and universal-property or
    tower/coarsening data;
  \item \Psix{}: record-event or check evidence together with the audit
    schema and its provenance and replay/check relations.
\end{itemize}
Each lift is selected and not unique, and forgetting from a higher level
loses structural content that is not recoverable without re-adding data.
\end{definition}

\begin{remark}[Global nonclaims for the role architecture]\label{rem:atlas-global-nonclaims}
The role architecture preserves the following nonclaims, all carried
forward into the admissibility regime of
Section~\ref{sec:admissibility-meta-theory} and the theorem chain of
Sections~\ref{sec:lower-bounds}--\ref{sec:decomposition-exhaustion}:
\begin{itemize}[topsep=2pt,itemsep=1pt]
  \item no full six-primitives algebra is claimed;
  \item no global primitive-independence theorem is claimed;
  \item no empirical realization is claimed;
  \item no uniqueness of primitive presentations is claimed;
  \item no uniqueness of selected lifts is claimed;
  \item no lower-level-to-higher-level determinacy is claimed;
  \item no instrument-free formulation is claimed;
  \item direct cross-level comparison remains inadmissible without
    translation.
\end{itemize}
\end{remark} \section{Admissibility and meta-theory}\label{sec:admissibility-meta-theory}



The admissibility regime is honest bookkeeping, elevated from a
per-description condition on \(D \in \FATCD{}\) to a meta-theoretic
discipline governing which claims about admissible descriptions enter
the theorem chain. Under honest bookkeeping, the flat
primitive-independence question is replaced by typed non-collapse,
level transitions are controlled by forgetting and selected lifts,
thresholds govern activation, instruments govern visibility, and
empirical-bridge claims are admitted only as threshold-dependent,
level-indexed, instrument-visible assertions. The admissibility
regime, recorded in the definitions and propositions of this section,
fixes these requirements and preserves the global nonclaims of
Remark~\ref{rem:atlas-global-nonclaims}.

\subsection{Honest bookkeeping as an admissibility regime}\label{subsec:honest-bookkeeping}

The audit condition of Definition~\ref{def:audited} records enough
data for each claim, promotion, translation, lift, forgetting,
visibility, and empirical-bridge assertion to be interrogated on its
own terms. The following definition stipulates that these record
requirements jointly constitute the admissibility regime, in a spirit
analogous to the axiomatic specification discipline of Hoare
logic~\citep{Hoare1969}.

\begin{definition}[Admissible claim under honest bookkeeping]\label{prop:honest-bookkeeping}
Honest bookkeeping is the admissibility regime for claims about a
description \(D \in \FATCD{}\). A claim that enters the theorem chain is
declared \emph{admissible} only if every assertion below carries an
explicit record in \(\Audit{D}\):
\begin{itemize}[topsep=2pt,itemsep=1pt]
  \item the claim itself, with its statement, hypotheses, and scope;
  \item every promotion, with its source level, target level, added data,
    losses, and preserved nonclaims;
  \item every translation, with its source context, target context,
    preserved data, suppressed data, and claim strength;
  \item every selected lift, marked selected and non-unique;
  \item every forgetting map, with its record of losses;
  \item every visibility claim, with both its visible and its suppressed
    data;
  \item every empirical-bridge assertion, with its substrate, observable,
    instrument, level map, threshold-witness relation, and
    suppressed-structure nonclaims.
\end{itemize}
Accordingly, unrecorded claims, implicit promotions, silent translations,
asserted-but-unselected lifts, lossless forgetting without proof,
visibility overreads, and bare empirical-realization assertions are all
inadmissible.
\end{definition}

\begin{remark}[Why the regime is necessary]\label{rem:why-honest-bookkeeping}
The bookkeeping requirement is a structural check, not a stylistic
one. Without explicit records, passive observation can be misread as
\Psix{}, mere packaging as \Pone{} descent, and route failure as
\Pthree{} mismatch. The regime prevents those misreadings from
entering the theorem chain.
\end{remark}


\subsection{Typed non-collapse}\label{subsec:typed-non-collapse}

The regime does not include flat primitive independence --- the claim
that \Pone{}, \ldots, \Psix{} are independent generators of a global
algebra. Its meta-theoretic counterpart here is typed non-collapse.

\begin{proposition}[Typed non-collapse]\label{prop:typed-non-collapse}
Under honest bookkeeping (Definition~\ref{prop:honest-bookkeeping}) and
the repaired alignments of Remark~\ref{rem:repaired-alignments}, for
every pair of primitive roles \((P_i, P_j)\) with \(i \ne j\), the
``forbidden collapse'' of the corresponding row in
Table~\ref{tab:boundary-matrix} is not admissible as an identification
of \(P_i\) with \(P_j\); the ``permitted interaction'' of the same row is
not thereby excluded and may remain admissible when its own bookkeeping
conditions are met.\footnote{Tracked in the formalization harness as a conditional schema (\texttt{SixBirds.typed\_non\_collapse}); see Appendix~\ref{app:mechanization}.}
\end{proposition}

\begin{proof}
Fix a pair \((P_i, P_j)\) with \(i \ne j\). The corresponding row of
Table~\ref{tab:boundary-matrix} records a ``forbidden collapse'': an
identification asserting that the \(P_i\)-content and the
\(P_j\)-content of a description are the same closure role, so that one
projection may be read off the other without further data. To be
admissible under Definition~\ref{prop:honest-bookkeeping}, such an
identification would have to enter the theorem chain as a claim with a
recorded statement and scope. We show that no such record can be
honest.

The two roles have distinct meanings fixed in
Definition~\ref{def:primitive-roles}: each of \Pone{}, \ldots, \Psix{}
carries a different closure function (descent obstruction,
macro-specification representability, enriched route mismatch,
stage/refinement structure, quotient packaging, and
audit/provenance/replay, respectively). An identification of \(P_i\)
with \(P_j\) would therefore equate two roles whose meanings differ.
Proposition~\ref{prop:boundary-distinctions} of
Section~\ref{sec:primitive-role-architecture} establishes that, under
the repaired alignments of Remark~\ref{rem:repaired-alignments}
together with the honest bookkeeping required by a \FATCD{}, the
forbidden collapse of every pair listed in
Table~\ref{tab:boundary-matrix} is ruled out as a locally admissible
identification of those two roles. In particular, the repaired
alignments forbid exactly the substitutions on which a collapse would
rely --- reading generic transition data as \Pthree{}, generic
path or derivation data as \Pfive{}, or generic semantic data as
\Pone{} --- so the cross-role record demanded by the collapse cannot be
supplied without misalignment. Hence the collapse fails the audit
condition of Definition~\ref{prop:honest-bookkeeping} and is
inadmissible.

It remains to check that ruling out the collapse does not also rule out
the row's ``permitted interaction''. The permitted interaction does not
identify \(P_i\) with \(P_j\); it records a legitimate joint occurrence
of the two distinct roles in a single description --- for instance,
\Pone{} descent audited by \Psix{}, or \Pthree{} route terms indexed by
\Pfour{} stages, as in Remark~\ref{rem:not-global-independence}. Such
an interaction carries its own honest-bookkeeping record for each role
separately and asserts no identification. Nothing in the argument above
touches it: the inadmissibility derived for the collapse rests solely
on the equation \(P_i = P_j\), which the permitted interaction never
asserts. Therefore the permitted interaction remains admissible
whenever its own bookkeeping conditions are met.
\end{proof}

\begin{remark}[Typed non-collapse versus global independence]\label{rem:typed-non-collapse-vs-independence}
Typed non-collapse is strictly weaker than global primitive
independence. It does not assert that the six roles are logically or
categorically independent in a global algebra; it asserts only that
the pairwise identifications listed as forbidden collapses in
Table~\ref{tab:boundary-matrix} are inadmissible under honest
bookkeeping, leaving the corresponding typed interactions available.
\end{remark}


\subsection{Level-profile architecture}\label{subsec:level-profile-architecture}

Typed non-collapse fixes which identifications are inadmissible; the
regime must also record the level at which each admissible claim is
made. The level trichotomy of Definition~\ref{def:level-trichotomy}
is therefore lifted from a role-by-role template to a regime-level
structure.

\begin{definition}[Level profile of an admissible claim]\label{prop:level-profile}
Every admissible claim about a description \(D \in \FATCD{}\) carries a
\emph{level profile}: each claim is located on the behavioral,
verification, or structural-categorical level, and each cross-level
assertion is accompanied by a translation record in the sense of
Definition~\ref{prop:honest-bookkeeping}. Direct comparison of claims at
different levels without a translation record is inadmissible.
\end{definition}

\subsection{Forgetting and selected lifts}\label{subsec:forgetting-lifts}

Level transitions are controlled by forgetting and selected lifts.

\begin{definition}[Admissible forgetting and selected lifts]\label{prop:forgetting-lifts}
Forgetting from the structural-categorical level to verification, and
from verification to behavioral, is admissible and is always lossy in
principle: an admissible forgetting carries a record of the losses in
\(\Audit{D}\). Lifting from a lower level to a higher level is
admissible only when accompanied by
\begin{itemize}[topsep=2pt,itemsep=1pt]
  \item the added data required to support the higher-level statement;
  \item the mark that the lift is selected and non-unique;
  \item explicit preservation of lower-level invariants and explicit
    losses where the lift is partial.
\end{itemize}
Lossless forgetting and unique lifts are inadmissible unless the losslessness
or uniqueness is itself proved and recorded.
\end{definition}

\begin{remark}[Lower-to-higher non-determinacy]\label{rem:lower-to-higher-non-determinacy}
Definition~\ref{prop:forgetting-lifts} formalizes the non-determinacy of
lower-level data: given behavioral or verification content alone, no
choice of structural-categorical content is forced by the admissibility
regime, and any asserted choice must be marked as a selected lift.
\end{remark}

\subsection{Activation thresholds}\label{subsec:activation-thresholds}

The regime must also state when a projected role is active. The
attachment protocol of Definition~\ref{def:threshold-attachment} is
lifted to a regime-level principle.

\begin{remark}[Activation thresholds]\label{prop:activation-thresholds}
At the regime level, the threshold-attachment discipline of
Definition~\ref{def:threshold-attachment} reads as follows. An
admissible active derived role projection \(\pi_i(D)\) exists only when
\(D\) carries the threshold, witness, level, collapse, and audit data
required by Definition~\ref{def:threshold-attachment}. Below-threshold
clusters carry an explicit \(\belowthr\) status, and boundary-collapsed
clusters carry \(\collapsedbc\). Arbitrary enrichment, adding structure
solely to force a role projection, is inadmissible unless the additions
are themselves honestly bookkept as raw features or annotations. This
is a restatement of the attachment protocol at regime level and adds no
content beyond Definition~\ref{def:threshold-attachment}.
\end{remark}

\subsection{Instrument-relative visibility and suppression}\label{subsec:visibility-suppression}

Threshold-governed activation still leaves open what an instrument
makes visible or suppresses. Admissibility is therefore
instrument-relative.

\begin{definition}[Admissible visibility record]\label{prop:visibility}
Every admissible claim carries an explicit \emph{visibility record}
comprising its visible content, its suppressed content, the instruments
under which the visibility and suppression are asserted, and the
resulting instrument-relative claim strength. Overreading a suppressed
claim as visible under a different instrument is inadmissible. Partial
visibility, in which a claim is visible only up to some coarsening or
translation, is admitted when the visibility record itself makes the
partiality explicit.
\end{definition}

\begin{remark}[No instrument-free truth]\label{rem:no-instrument-free}
Definition~\ref{prop:visibility} does not claim that instruments are
free of distortion or that any single instrument is canonical. It
asserts only that admissible claims carry an instrument-relative
visibility record. Instrument completeness is deferred to
Section~\ref{sec:future-work}.
\end{remark}

\subsection{Empirical-bridge admissibility}\label{subsec:empirical-bridge}

Any claim crossing from the closure-algebraic side of the framework
to a substrate-level observable must pass through the
empirical-bridge discipline, related to but narrower than
the empirical-adequacy view of theories~\citep{vanFraassen1980}.

\begin{definition}[Admissible empirical-bridge claim]\label{prop:empirical-bridge}
An empirical-bridge claim about \(D \in \FATCD{}\) is admissible only when
it carries, in \(\Audit{D}\),
\begin{itemize}[topsep=2pt,itemsep=1pt]
  \item the substrate and observable to which the bridge relates;
  \item the instrument under which the observable is measured;
  \item the level map between the closure-algebraic content and the
    substrate-level content;
  \item a threshold-witness relation recording when the bridge is
    crossed;
  \item the nonclaims suppressed by the bridge record, in particular any
    structural content of \(D\) that is not tracked by the substrate or
    instrument.
\end{itemize}
Empirical realization --- the claim that \(D\) is itself a substrate
observation --- is stronger than empirical-bridge admissibility and is
not asserted by any empirical-bridge claim.
\end{definition}


\subsection{Nonclaim closure across the admissibility regime}\label{subsec:nonclaim-closure}

The admissibility regime, recorded in the definitions and propositions
of this section --- honest bookkeeping
(Definition~\ref{prop:honest-bookkeeping}), typed non-collapse
(Proposition~\ref{prop:typed-non-collapse}), the level-profile
architecture (Definition~\ref{prop:level-profile}), admissible
forgetting and selected lifts (Definition~\ref{prop:forgetting-lifts}),
activation thresholds (Remark~\ref{prop:activation-thresholds}),
admissible visibility (Definition~\ref{prop:visibility}), and
admissible empirical-bridge claims (Definition~\ref{prop:empirical-bridge})
--- together preserve the global nonclaims of
Remark~\ref{rem:atlas-global-nonclaims} and block a corresponding
class of overclaims.

\begin{proposition}[Admissibility-regime nonclaim closure]\label{prop:nonclaim-closure}
Under the admissibility regime, the following remain non-claims:
\begin{itemize}[topsep=2pt,itemsep=1pt]
  \item no full six-primitives algebra is claimed;
  \item no global primitive-independence theorem is claimed;
  \item no empirical realization is claimed;
  \item no unique lift theorem is claimed;
  \item no lower-level-to-higher-level determinacy is claimed;
  \item no instrument-free truth claim is made;
  \item direct cross-level comparison remains inadmissible without
    translation;
  \item forgetting remains lossy unless losslessness is explicitly proved;
  \item empirical-bridge claims require threshold data, level maps,
    instrument visibility records, and suppressed-structure nonclaims.
\end{itemize}\footnote{Tracked in the formalization harness as a conditional schema (\texttt{SixBirds.nonclaim\_closure}); see Appendix~\ref{app:mechanization}.}
\end{proposition}

\begin{proof}
Each requirement of the regime blocks a specific class of overclaims,
and the listed non-claims are exactly the residue.

Honest bookkeeping (Definition~\ref{prop:honest-bookkeeping}) declares
inadmissible every unrecorded claim, implicit promotion, silent
translation, asserted-but-unselected lift, lossless forgetting without
proof, visibility overread, and bare empirical-realization assertion;
this is the audit condition of Definition~\ref{def:audited} read as an
admissibility filter.

Typed non-collapse (Proposition~\ref{prop:typed-non-collapse}) blocks
the forbidden primitive collapses of Table~\ref{tab:boundary-matrix}
and, with them, the flat-independence overclaim that would treat the
six roles as independent generators of a global algebra; by the same
proposition it leaves the restricted typed boundary interactions
available, so the standing non-claim is the global
primitive-independence theorem and the full six-primitives algebra, not
the absence of interaction.

Level-profile architecture (Definition~\ref{prop:level-profile}) blocks
direct cross-level comparison without a translation record, which is
exactly the non-claim that cross-level comparison remains inadmissible
without translation.

Admissible forgetting and selected lifts
(Definition~\ref{prop:forgetting-lifts}) block lossless forgetting,
unique lifts, and lower-to-higher determinacy: forgetting is lossy
unless losslessness is proved and recorded, no lift is admitted as
unique, and, by
Remark~\ref{rem:lower-to-higher-non-determinacy}, no structural-categorical
content is forced by behavioral or verification content alone. These are
the non-claims that forgetting remains lossy, that no unique-lift
theorem is claimed, and that no lower-to-higher determinacy is claimed.

Activation thresholds (Remark~\ref{prop:activation-thresholds}) block
arbitrary enrichment and threshold-free nontriviality claims, since an
active projection is admissible only with the threshold, witness, level,
collapse, and audit data of Definition~\ref{def:threshold-attachment}.

Instrument-relative visibility (Definition~\ref{prop:visibility}) blocks
instrument-free truth claims and visible-to-suppressed overreads, which
yields the non-claim that no instrument-free truth claim is made.

Empirical-bridge admissibility (Definition~\ref{prop:empirical-bridge})
blocks empirical-realization overclaims and any bridge claim lacking
threshold data, a level map, an instrument visibility record, or
suppressed-structure nonclaims; this yields the non-claim that no
empirical realization is claimed and that empirical-bridge claims carry
the required data.

Collecting these residues gives exactly the listed non-claims. Those that
restate a global nonclaim of Remark~\ref{rem:atlas-global-nonclaims} are
thereby carried forward into the regime; the two further items --- that
forgetting remains lossy absent a recorded proof of losslessness, and that
empirical-bridge claims carry their threshold, level-map, visibility, and
suppressed-structure data --- are additional disciplines the regime imposes
at this level.
\end{proof}
 \section{The FATCD domain and the exact-six question}\label{sec:fatcd-domain}



The scoped exact-six theorem quantifies over a specific class of
admissible descriptions; this section defines that class. A
\emph{finite audited typed closure description}, abbreviated \FATCD{},
is a finitely presented collection of typed raw records and
annotations subject to finiteness, closure, and honest-bookkeeping
conditions, sitting in the broader finite-structure tradition of
finite model theory~\citep{Libkin2004}. Threshold annotations attach to raw features before any
role projection is attempted; the active projections
\(\pi_1(D),\dots,\pi_6(D)\) for the primitives \Pone{}--\Psix{} are
derived interpretations of raw material, not conditions of domain
membership. Candidate raw features not absorbed into such projections
are carried into an explicit residual classification, so seventh-role
threats are not excluded by definition. The section closes with
non-circularity and non-overbreadth schemas and formulates the
exact-six question precisely over \FATCD{}: whether every
\(D \in \FATCD{}\) admits a decomposition/exhaustion record with
exactly six role channels. The theorem chain answering that question
appears in
Sections~\ref{sec:lower-bounds}--\ref{sec:decomposition-exhaustion}.

The domain is framed in the finite theory-package setting of
Foundations~I~\citep{Tsiokos2026SixBirdsFoundations}, where one works
with tuples \(\mathcal{T} = (Z, f, \Sigma_f, E, \mathcal{A})\)
carrying a lens, a completion operator, and an audit functional. \FATCD{} does not
identify a description with such a tuple; it extracts from that
setting the structural requirements relevant to the scoped exact-six
question: finite typed raw data, closure under references, sources,
targets, endpoints, and annotations, and honest bookkeeping in the
sense of Section~\ref{sec:admissibility-meta-theory}. Optimization,
objective, and richer semantic material are not excluded at the
outset; they belong in the domain when finitely represented and
honestly recorded as raw or annotated data.

\subsection{Raw grammar}\label{subsec:raw-grammar}

We begin with the neutral typed raw grammar from which every element
of \FATCD{} is built. The alphabet is fixed in advance of the role
vocabulary and follows standard typed-record
conventions~\citep{CardelliMitchell1991}.

\begin{definition}[Raw record]\label{def:raw-record}
A \emph{raw record} is a finite typed syntactic or relational feature of a
description. The admissible record kinds are
\begin{itemize}[topsep=2pt,itemsep=1pt]
  \item \emph{object} and \emph{state} records;
  \item \emph{relation} records, \emph{map} records, and \emph{partial-map}
    records;
  \item \emph{predicate} records and \emph{comparison} records;
  \item \emph{path} or \emph{derivation} records;
  \item \emph{index}, \emph{order}, \emph{transition}, and \emph{refinement}
    records;
  \item \emph{trace}, \emph{event}, and \emph{provenance} records;
  \item \emph{replay} and \emph{check} records.
\end{itemize}
A raw record is not, by itself, a \Pone{}--\Psix{} role.
\end{definition}

\begin{definition}[Annotations]\label{def:annotations}
In addition to raw records, a description admits \emph{annotations} of six
declared kinds:
\begin{itemize}[topsep=2pt,itemsep=1pt]
  \item level annotations \(\mathrm{Ann}_L\);
  \item threshold annotations \(\mathrm{Ann}_\Theta\);
  \item lift annotations \(\mathrm{Ann}_{\mathrm{Lift}}\);
  \item forgetting annotations \(\mathrm{Ann}_{\mathrm{Forget}}\);
  \item instrument-visibility annotations \(\mathrm{Ann}_{\mathrm{Vis}}\);
  \item empirical-bridge annotations \(\mathrm{Ann}_{\mathrm{Bridge}}\).
\end{itemize}
These annotation kinds belong to the raw grammar; the finiteness and closure
conditions on how they attach to declared records are imposed in the next
subsection. For a description \(D\), we write \(\Raw{D}\) for the finite
collection of raw records and annotations of \(D\).
\end{definition}

\begin{remark}[Independence of the feature alphabet]\label{rem:alphabet-independence}
The record and annotation kinds in Definitions~\ref{def:raw-record} and
\ref{def:annotations} are syntactic or relational entities. No kind is named
\Pone{}, \ldots, \Psix{}. Consequently, membership in \FATCD{} does not
require any derived role projection, does not require an exact-six
decomposition/exhaustion record, and does not exclude candidate seventh-role
features by definition. The role vocabulary will enter only in
\S\ref{subsec:derived-role-projections} as a derived interpretation of raw
features under threshold and witness data.
\end{remark}

\subsection{Finiteness, closure, and audit conditions}\label{subsec:conditions}

Three further restrictions turn the raw grammar into a domain of
admissible descriptions: finiteness, closure, and audit.

\begin{definition}[Finite description]\label{def:finite}
A description \(D\) is \emph{finite} when
\begin{itemize}[topsep=2pt,itemsep=1pt]
  \item \(\Raw{D}\) contains finitely many records and annotations;
  \item every record or annotation carries finite typed fields;
  \item every relation, map, path, trace, or annotation is finitely
    represented;
  \item every reference in a record points to a declared finite record of
    \(D\);
  \item no infinite semantic, probabilistic, topological, geometric, causal,
    optimization, or resource structure is admitted unless it is finitely
    represented and bookkept as a raw feature of \(D\).
\end{itemize}
\end{definition}

\begin{definition}[Closed description]\label{def:closed}
A finite description \(D\) satisfies the \emph{closure condition} \(\Closed{D}\)
when
\begin{itemize}[topsep=2pt,itemsep=1pt]
  \item every source and target of every relation or map record is a declared
    record of \(D\);
  \item every endpoint of every path or derivation record is a declared record
    of \(D\);
  \item every index, transition, order, or refinement reference is a declared
    record of \(D\);
  \item every trace, event, provenance, replay, or check record references
    declared records of \(D\);
  \item every level, threshold, lift, forgetting, visibility, or
    empirical-bridge annotation references declared records of \(D\);
  \item every claim carried by \(D\) has a bookkeeping record.
\end{itemize}
\end{definition}

\begin{definition}[Audited description]\label{def:audited}
A description \(D\) is \emph{audited} when every claim carried by \(D\) has an
honest bookkeeping record. The collection of such records is denoted
\(\Audit{D} \subseteq \Raw{D}\), subject to the following discipline:
\begin{itemize}[topsep=2pt,itemsep=1pt]
  \item every promotion records its source level, target level, added data,
    losses, and nonclaims;
  \item every translation records its source context, target context,
    preserved data, suppressed data, and claim strength;
  \item every lift is marked as selected, not unique;
  \item every forgetting map records its losses;
  \item every visibility claim records both visible and suppressed data;
  \item every empirical-bridge claim records substrate, observable,
    instrument, level map, threshold witness relation, and
    suppressed-structure nonclaims.
\end{itemize}
\end{definition}

\subsection{Finite audited typed closure descriptions}\label{subsec:fatcd}

Combining the grammar with the three admissibility conditions yields
the domain over which the exact-six question is posed.

\begin{definition}[Finite audited typed closure description]\label{def:fatcd}
A \emph{finite audited typed closure description} is a finitely presented,
closed, audited collection of typed raw records and annotations, equipped with
honest bookkeeping for claims, levels, thresholds, visibility, forgetting,
lifts, and empirical-bridge assertions where present. Equivalently, it is a
description \(D\) whose raw content \(\Raw{D}\) is drawn from the kinds
declared in Definitions~\ref{def:raw-record} and~\ref{def:annotations}, such
that \(D\) is finite in the sense of Definition~\ref{def:finite}, satisfies
the closure condition \(\Closed{D}\) of Definition~\ref{def:closed}, and is
audited in the sense of Definition~\ref{def:audited}. We write \FATCD{} for
the class of all such descriptions, and \(D \in \FATCD{}\) for membership.
\end{definition}

\begin{remark}[No minimum richness]\label{rem:fatcd-minimality}
\FATCD{} imposes no lower bound on description size. A description
consisting of a single object record together with the bookkeeping
required by its claims may be admissible, as may a richer description
carrying transitions, paths or derivations, annotations of every kind,
and empirical-bridge annotations. The domain is defined by structural
discipline, not by the presence of any particular feature kind.
\end{remark}

\subsection{Threshold attachment}\label{subsec:threshold-attachment}

Membership in \FATCD{} is fixed; the next step is to specify how
role-theoretic claims may attach to the raw data of an admissible
description. Role projections in \FATCD{} enter only through a
disciplined attachment of threshold and witness data to raw features.
The condition below is the structural precondition for the role
vocabulary introduced in the next subsection.

\begin{definition}[Threshold attachment]\label{def:threshold-attachment}
A threshold annotation \(\mathrm{Ann}_\Theta \in \Raw{D}\) attaches first to
a finite cluster \(C \subseteq \Raw{D}\) of raw features of \(D\), where \(C\)
may consist of a single raw feature. No role projection is attempted on \(C\)
until all of the following are present in \(\Raw{D}\):
\begin{itemize}[topsep=2pt,itemsep=1pt]
  \item the raw features of \(C\) themselves;
  \item a threshold annotation \(\mathrm{Ann}_\Theta\) attached to \(C\);
  \item a witness schema specifying what in \(C\) would support the attempted
    projection;
  \item a record of the collapse cases that would invalidate the witness;
  \item a level annotation \(\mathrm{Ann}_L\) attached to \(C\);
  \item an audit record in \(\Audit{D}\) honestly bookkeeping the claim.
\end{itemize}
\end{definition}

\begin{remark}[Pre-projection discipline]\label{rem:pre-projection}
The items of Definition~\ref{def:threshold-attachment} form a structural
precondition, not a process narrative. What matters is that each item is
present explicitly in \(\Raw{D}\) before a role projection is asserted, so
that no threshold, witness, collapse case, or level is supplied tacitly by
the act of projection itself.
\end{remark}

\subsection{Derived role projections}\label{subsec:derived-role-projections}

The role vocabulary \Pone{}--\Psix{} now enters \FATCD{}, not as a
constraint on membership but as a derived interpretation of raw
features conditioned on the attachment data of
\S\ref{subsec:threshold-attachment}.

\begin{definition}[Derived role projection]\label{def:role-projection}
A \emph{derived role projection} for a description \(D \in \FATCD{}\) at
role index \(i \in \{1,\ldots,6\}\) is an active derived interpretation
\(\pi_i(D)\) of raw features of \(D\) as the role \(\mathrm{P}_i\), supported
by threshold data, witness schema, level assignment, and honest bookkeeping.
Concretely, \(\pi_i(D)\) is specified by
\begin{itemize}[topsep=2pt,itemsep=1pt]
  \item the proposed role name, one of \Pone{}, \Ptwo{}, \Pthree{},
    \Pfour{}, \Pfive{}, \Psix{};
  \item a raw feature cluster \(C \subseteq \Raw{D}\) on which the role is
    interpreted;
  \item the threshold data \(\Theta_i\) attached to \(C\) in the sense of
    Definition~\ref{def:threshold-attachment};
  \item a witness schema \(W_i\) specifying what constitutes a witness for
    the candidate role under its fixed meaning;
  \item a level assignment \(\mathrm{Level}_i\);
  \item an evidence or certificate type;
  \item a record of collapse cases avoided or accepted;
  \item loss and lift annotations for any level crossings involved;
  \item the explicit nonclaims preserved by the projection.
\end{itemize}
The supporting data for \(\pi_i(D)\) are carried by \(\Raw{D}\) and certified
by \(\Audit{D}\) in the sense of Definition~\ref{def:audited}.
\end{definition}

\begin{remark}[Fixed role meanings]\label{rem:fixed-role-meanings}
The role names \Pone{}--\Psix{} are defined by their meanings in
Section~\ref{sec:primitive-role-architecture}: \Pone{} concerns
update-vs-quotient compatibility and descent obstruction; \Ptwo{} concerns
macro-specification representability; \Pthree{} concerns enriched
route and coherence mismatch, not a directionality certificate; \Pfour{}
concerns stage, order, transition, and refinement structure; \Pfive{}
concerns quotient packaging; and
\Psix{} concerns audit, bookkeeping, provenance, replay, and checkability.
The witness schema \(W_i\) in Definition~\ref{def:role-projection} must be
aligned with these meanings; a cluster whose raw features do not support the
meaning of the proposed role contributes no \(\pi_i(D)\).
\end{remark}

\begin{remark}[Projection and membership are independent]\label{rem:projection-vs-membership}
Whether a description \(D\) supports a given derived role projection
\(\pi_i(D)\) is not part of the admissibility conditions of
Definitions~\ref{def:finite}--\ref{def:audited}. A description
\(D \in \FATCD{}\) may carry no derived role projection, projections
for some roles, or projections for all six; its membership in
\FATCD{} is unaffected. The decomposition/exhaustion record of
Section~\ref{sec:decomposition-exhaustion} accordingly assigns a
status to every role channel of \(D\), only some of which need be
active.
\end{remark}

\subsection{Residual classification scheme}\label{subsec:residual-classification}

Some raw features remain role-theoretically unsettled at the
raw-grammar level. \FATCD{} provides an explicit scheme for
classifying them.

\begin{definition}[Residual feature; residual classification scheme]\label{def:residual-scheme}
Let \(D \in \FATCD{}\). A \emph{residual feature} of \(D\) is a raw record
or annotation in \(\Raw{D}\) whose role-theoretic disposition is not fixed
solely by its raw-grammar kind. The \emph{residual classification scheme}
assigns to each residual feature \(r\) a single label
\(\operatorname{Classify}(r)\) drawn from the list
\begin{enumerate}[label=(\roman*),topsep=2pt,itemsep=1pt]
  \item already absorbed into a derived \Pone{}--\Psix{} projection
    \(\pi_i(D)\);
  \item composite of such projections;
  \item presentation variant of an already classified feature;
  \item level shift of an already classified feature;
  \item activation-threshold refinement;
  \item instrument-visibility artifact;
  \item empirical-bridge annotation;
  \item outside-domain structure;
  \item already-covered bookkeeping or annotation data;
  \item unresolved candidate seventh-role threat;
  \item genuine seventh irreducible typed role.
\end{enumerate}
The scheme is declarative. At the domain-definition level, \FATCD{} records
the available labels but neither proves that every residual feature lands in
categories (i)--(ix) nor proves that categories (x) and (xi) are empty. Those
claims are deferred to Section~\ref{sec:upper-bound-classification}.
\end{definition}

\begin{remark}[Optimization and objective data are not excluded]\label{rem:optimization-not-excluded}
Raw features expressing optimization, objective, or selection criteria
pass through the scheme of Definition~\ref{def:residual-scheme}; they
are not dismissed by definition. Under the fixed role meanings
(Remark~\ref{rem:fixed-role-meanings}), typical classifications
include: a predicate or specification record possibly supporting a
\Ptwo{} projection; a map or transition-selection record possibly
supporting a \Pone{} or \Pfour{} projection; a route-comparison or
route-choice record possibly supporting a \Pthree{} projection; an
audit or check record possibly supporting a \Psix{} projection; an
empirical-bridge annotation; or an outside-domain structure when the
feature requires non-finite or unbookkept semantics. A feature that
fits none of these is recorded under category~(x).
\end{remark}

\subsection{Non-circularity and non-overbreadth}\label{subsec:non-circularity-non-overbreadth}

Two proposition schemas record the definition-theoretic facts used by
the later theorem chain: the domain does not presuppose the answer to
the exact-six question, and it is structurally restrictive enough
for that question to remain substantive.

\begin{proposition}[Non-circularity schema for \FATCD{}]\label{prop:non-circularity}
Membership in \FATCD{} is not defined by requiring a
decomposition/exhaustion record or by requiring any derived role
projection to exist.\footnote{Tracked in the formalization harness as a conditional schema (\texttt{SixBirds.non\_circularity}); see Appendix~\ref{app:mechanization}.}
\end{proposition}

\begin{proof}
We inspect the clauses that govern membership, namely those of
Definitions~\ref{def:raw-record}--\ref{def:fatcd}, and the residual
classification scheme of Definition~\ref{def:residual-scheme}.
\begin{enumerate}[label=(\roman*),topsep=2pt,itemsep=1pt]
  \item No clause of Definitions~\ref{def:raw-record},
    \ref{def:annotations}, \ref{def:finite}, \ref{def:closed}, or
    \ref{def:audited} quantifies over the role index \(i\) or requires a
    role-channel assignment. The record and annotation kinds
    (Definitions~\ref{def:raw-record} and~\ref{def:annotations}) are
    syntactic or relational entities; the finiteness, closure, and audit
    conditions (Definitions~\ref{def:finite}--\ref{def:audited}) constrain
    cardinality, field types, reference resolution, and bookkeeping
    discipline. None of these conditions mentions \(\pi_i\), a role channel,
    or a decomposition/exhaustion record. Consequently membership cannot
    presuppose that any such projection or record exists.
  \item No membership clause excludes a candidate seventh-role feature by
    definition. A candidate seventh-role feature is, at the raw-grammar
    level, a record or annotation of one of the declared kinds; nothing in
    Definitions~\ref{def:finite}--\ref{def:audited} filters a feature out on
    the ground that it might resist absorption into a
    \Pone{}--\Psix{} projection. Such a feature is therefore admissible
    whenever it is finite, closed, and audited, and is carried forward to
    the residual classification scheme rather than barred at the door.
  \item The residual classification scheme
    (Definition~\ref{def:residual-scheme}) is declarative and carries the
    categories (x) and (xi)---an unresolved candidate seventh-role threat,
    and a genuine seventh irreducible typed role---as explicit failure
    options. At the domain-definition level the scheme records the available
    labels but does not assert that those two categories are empty, so it
    cannot smuggle the exact-six answer into the definition of the domain.
\end{enumerate}
Since none of (i)--(iii) ties admissibility to the existence of a role
projection or a six-channel record, membership in \FATCD{} is independent of
the exact-six question, which is what was claimed.
\end{proof}

\begin{proposition}[Non-overbreadth schema for \FATCD{}]\label{prop:non-overbreadth}
\FATCD{} is structurally restricted and does not admit arbitrary
syntactic material: arbitrary unaudited syntactic material cannot enter
\FATCD{}, and mentioning a role name without the supporting threshold,
witness, level, collapse, and bookkeeping data does not by itself produce a
derived role projection.\footnote{Tracked in the formalization harness as a conditional schema (\texttt{SixBirds.non\_overbreadth}); see Appendix~\ref{app:mechanization}. The Lean declaration is identical in statement and proof to \texttt{SixBirds.non\_circularity} and re-extracts the admissibility fields of the record rather than establishing the structural-restriction content of the proposition.}
\end{proposition}

\begin{proof}
Let \(D \in \FATCD{}\). We inspect the admissibility conditions of
Definitions~\ref{def:finite}--\ref{def:threshold-attachment} in turn.
\begin{enumerate}[label=(\roman*),topsep=2pt,itemsep=1pt]
  \item By Definition~\ref{def:finite}, \(\Raw{D}\) contains finitely many
    records and annotations, each record or annotation carries finite typed
    fields, and no infinite semantic, probabilistic, topological, geometric,
    causal, optimization, or resource structure is admitted unless finitely
    represented and bookkept as a raw feature. Thus unbounded or
    unrepresented material is excluded at the level of finiteness.
  \item By Definition~\ref{def:closed}, every source, target, endpoint,
    index, transition, order, refinement, trace, event, provenance, replay,
    check, and annotation reference is bound to a declared record of \(D\).
    Hence no dangling or externally defined reference can occur, and the
    syntactic material of \(D\) is self-contained.
  \item By Definition~\ref{def:audited}, every claim carried by \(D\) is
    bookkept, and every promotion, translation, lift, forgetting, visibility,
    and empirical-bridge assertion carries the structured audit record
    required there. Therefore an assertion lacking its honest bookkeeping
    record cannot be part of an admissible \(D\).
  \item By Definition~\ref{def:threshold-attachment}, a role projection
    cannot be asserted on a cluster \(C\) until the raw features of \(C\), a
    threshold annotation, a witness schema, a record of collapse cases, a
    level annotation, and an audit record are all present in \(\Raw{D}\).
    Hence a role name supplied without this attachment data does not, by the
    act of naming alone, produce a projection.
\end{enumerate}
Combining (i)--(iii), arbitrary unaudited syntactic material cannot enter
\FATCD{}, since every admissible \(D\) is finite, reference-closed, and
honestly bookkept. By (iv), mentioning a role name without the supporting
threshold, witness, level, collapse, and bookkeeping data does not by itself
produce a derived role projection. This establishes both halves of the
claim.
\end{proof}

\subsection{The scoped exact-six question over \FATCD{}}\label{subsec:scoped-exact-six-question}

We can now state the scoped exact-six question precisely.

\begin{definition}[Scoped exact-six question]\label{def:scoped-exact-six-question}
The \emph{scoped exact-six question} over \FATCD{} asks whether, for every
\(D \in \FATCD{}\), there exists at least one decomposition/exhaustion record
\(\Dec{D}\) whose role-channel structure consists of exactly the six typed
closure-role channels \Pone{}, \Ptwo{}, \Pthree{}, \Pfour{}, \Pfive{},
\Psix{}, each assigned exactly one of the five statuses later formalized in
Section~\ref{sec:decomposition-exhaustion}, with all residual features of
\(D\) classified under the scheme of Definition~\ref{def:residual-scheme}
without recourse to categories (x) or (xi). No uniqueness is built into this
question. The affirmative answer is denoted \(\ScopedExactSix{\FATCD{}}\), and is
established in Section~\ref{sec:scoped-exact-six} conditionally on the recognition
hypothesis introduced in Section~\ref{sec:upper-bound-classification}
(Definition~\ref{def:recognition-hypothesis}), with the three principal threat
classes settled unconditionally.
\end{definition}

\begin{remark}[Channels, not active projections]\label{rem:channel-vs-activation-reminder}
An affirmative answer to the scoped exact-six question asserts the
existence of six role channels per description, not six active derived
projections per description. The distinction introduced in
Remark~\ref{rem:projection-vs-membership} and formalized in
Section~\ref{sec:decomposition-exhaustion} is preserved throughout.
\end{remark}

\begin{remark}[Relation to the Foundations theorem]\label{rem:fatcd-foundations-relation}
Proving \(\ScopedExactSix{\FATCD{}}\) is not a re-derivation of the
canonical unavoidability result of Foundations~I~\citep{Tsiokos2026SixBirdsFoundations},
and \(\ScopedExactSix{\FATCD{}}\) does not follow from that result
alone. The Foundations theorem derives the \Pone{}--\Psix{} vocabulary
under minimal hypotheses of composability with limited interface
access; the scoped exact-six claim asserts that, within the finite
audited typed closure regime of \FATCD{} and under the recognition hypothesis of
Section~\ref{sec:upper-bound-classification}, every admissible description
admits a six-channel decomposition/exhaustion record without a seventh
irreducible typed role in the covered scope. The two results share
vocabulary but address different questions.
\end{remark}
 \section{Six-role lower bounds}\label{sec:lower-bounds}




The exact-six question of
\S\ref{subsec:scoped-exact-six-question} has two directions. The
lower-bound direction asserts that each of \Pone{}, \ldots, \Psix{}
is instantiated in \FATCD{}; the upper-bound classification of
Section~\ref{sec:upper-bound-classification} will assert, under a recognition
hypothesis, that no covered seventh role is needed. This section establishes the
lower-bound direction by constructing one admissible witness for each
role. The content is occurrence, not exhaustion: lower bounds alone
do not answer the exact-six question, but they rule out the
trivialization in which some role has no admissible witness in the
domain.

\subsection{Lower-bound theorem and common witness requirements}\label{subsec:lower-bound-theorem}

\begin{theorem}[Six-role lower bounds]\label{thm:six-role-lower-bounds}
For each role index \(i \in \{1,\ldots,6\}\), there exists a description
\(D_i \in \FATCD{}\) and a derived active projection
\(\pi_i(D_i)\) witnessing the role channel \(P_i\) in the sense of
Definition~\ref{def:role-projection}. In particular, each of
\Pone{}, \Ptwo{}, \Pthree{}, \Pfour{}, \Pfive{}, \Psix{} is instantiated
by at least one admissible description in \FATCD{}. This is an occurrence
theorem, not an exhaustion theorem.\footnote{Tracked in the formalization harness as a conditional schema (\texttt{SixBirds.six\_role\_lower\_bounds}); see Appendix~\ref{app:mechanization}.}
\end{theorem}

A witness for role \(P_i\) consists of the data specified below,
following the role-projection discipline of
\S\ref{subsec:derived-role-projections} and the threshold-attachment
discipline of \S\ref{subsec:threshold-attachment}.

\begin{definition}[Common witness data for role \(P_i\)]\label{def:common-witness-data}
A \emph{witness of role \(P_i\) in \FATCD{}} is a tuple
\[
\bigl(D_i,\; C_i,\; \Theta_i,\; W_i,\; \mathrm{Level}_i,\;
v_i^{\mathrm{fin}},\; v_i^{\mathrm{cl}},\; v_i^{\mathrm{aud}},\;
n_i,\; \pi_i(D_i)\bigr)
\]
where
\begin{itemize}[topsep=2pt,itemsep=1pt]
  \item \(D_i \in \FATCD{}\) is the ambient admissible description;
  \item \(C_i \subseteq \Raw{D_i}\) is the raw feature cluster on which
    the projection is interpreted;
  \item \(\Theta_i\) is the threshold annotation attached to \(C_i\);
  \item \(W_i\) is the witness schema for the role meaning of \(P_i\)
    fixed in Definition~\ref{def:primitive-roles};
  \item \(\mathrm{Level}_i\) is the level assignment in the sense of
    Definition~\ref{def:level-trichotomy};
  \item \(v_i^{\mathrm{fin}}, v_i^{\mathrm{cl}}, v_i^{\mathrm{aud}}\) are
    the finite-description, closure, and audit/bookkeeping verifications
    for \(D_i\);
  \item \(n_i\) is the explicit nonclaim record carried by the witness;
  \item \(\pi_i(D_i)\) is the derived role projection of
    Definition~\ref{def:role-projection}.
\end{itemize}
A witness is admissible only if every item in this tuple is present and
correctly aligned with the role meaning of \(P_i\).
\end{definition}

The proof of Theorem~\ref{thm:six-role-lower-bounds} proceeds role by
role: for each \(i \in \{1,\ldots,6\}\) we exhibit an admissible
tuple of the form in Definition~\ref{def:common-witness-data}. The
six constructions are collected in
\S\ref{subsec:per-role-witnesses}.

\subsection{Per-role witnesses}\label{subsec:per-role-witnesses}

Each construction below specifies a raw feature cluster, threshold
annotation, witness schema, level assignment, derived active
projection, and explicit nonclaim record. The finite-description,
closure, and audit/bookkeeping verifications are routine once these
data are placed in a description built from the raw grammar of
\S\ref{subsec:raw-grammar} and checked against the admissibility
conditions of \S\ref{subsec:conditions}.

\begin{proposition}[\Pone{} lower bound]\label{prop:p1-witness}
There exists \(D_1 \in \FATCD{}\) with a derived active projection
\(\pi_1(D_1)\) witnessing \Pone{}.\footnote{The Lean development establishes this for one concrete description \texttt{witnessDescription Role.P1} under parametric hypotheses (alignment, threshold attachment, and finite/closed/audited verifications of that fixed witness); the narrowing is recorded in Appendix~\ref{app:mechanization} (\texttt{SixBirds.p1\_witness}).}
\end{proposition}

\begin{proof}
Let \(D_1\) carry object records \(x, y\), a relation record
\(x \sim y\), a selected update map record \(F\), and a comparison
record exhibiting \(F(x) \not\sim' F(y)\) in the target relation
\(\sim'\). Attach a threshold annotation
\(\Theta_1 = \mathrm{Ann}_\Theta[\text{descent obstruction}]\) to the
cluster \(C_1\) comprising these records, and record the witness schema
\(W_1\) that detects update-vs-quotient incompatibility between \(\sim\)
and \(\sim'\) under \(F\). Take \(\mathrm{Level}_1\) on the verification
level, and let \(\pi_1(D_1)\) be the derived active projection
interpreting \(C_1\) via \((\Theta_1,W_1,\mathrm{Level}_1)\). The
finiteness, closure, and audit verifications \(v_1^{\mathrm{fin}},
v_1^{\mathrm{cl}}, v_1^{\mathrm{aud}}\) then attach routinely. The
nonclaim record \(n_1\) excludes any \Pfive{}-packaging
identification (Remark~\ref{rem:repaired-alignments}) and any claim of
an unconditional descended map in the obstruction case.
\end{proof}

\begin{proposition}[\Ptwo{} lower bound]\label{prop:p2-witness}
There exists \(D_2 \in \FATCD{}\) with a derived active projection
\(\pi_2(D_2)\) witnessing \Ptwo{}.\footnote{The Lean development establishes this for one concrete description \texttt{witnessDescription Role.P2} under parametric hypotheses (alignment, threshold attachment, and finite/closed/audited verifications of that fixed witness); the narrowing is recorded in Appendix~\ref{app:mechanization} (\texttt{SixBirds.p2\_witness}).}
\end{proposition}

\begin{proof}
Let \(D_2\) carry object records together with a predicate record \(s\)
playing the role of a macro specification, a relation record
\(\mathrm{Real}\) connecting realizers to \(s\), and a comparison
record showing that the realizer relation is not reducible to
quotient-class equality alone. Attach
\(\Theta_2 = \mathrm{Ann}_\Theta[\text{nontrivial representability}]\)
to the resulting cluster \(C_2\), and record \(W_2\) as the schema
checking that \(s\) has or lacks a realizer through \(\mathrm{Real}\) in
a way not reducible to quotient-class equality alone. Take
\(\mathrm{Level}_2\) on the verification level, and let \(\pi_2(D_2)\)
be the derived active projection determined by
\((C_2,\Theta_2,W_2,\mathrm{Level}_2)\). The finiteness, closure, and
audit checks then follow from the construction. The
nonclaim \(n_2\) excludes collapse to \Pfive{} quotient equality alone
and excludes arbitrary relation promotion.
\end{proof}

\begin{proposition}[\Pthree{} lower bound]\label{prop:p3-witness}
There exists \(D_3 \in \FATCD{}\) with a derived active projection
\(\pi_3(D_3)\) witnessing \Pthree{}.\footnote{The Lean development establishes this for one concrete description \texttt{witnessDescription Role.P3} under parametric hypotheses (alignment, threshold attachment, and finite/closed/audited verifications of that fixed witness); the narrowing is recorded in Appendix~\ref{app:mechanization} (\texttt{SixBirds.p3\_witness}).}
\end{proposition}

\begin{proof}
Let \(D_3\) carry object records \(u, v\), two path records
\(\rho_1, \rho_2\) both from \(u\) to \(v\) generated by the enriched
route grammar, and a comparison record recording that \(\rho_1\) and
\(\rho_2\) are distinguishable by a non-reducibility certificate.
Attach \(\Theta_3 = \mathrm{Ann}_\Theta[\text{route mismatch}]\) to the
cluster \(C_3 = \{u, v, \rho_1, \rho_2, \text{comparison}\}\), and
record \(W_3\) as the schema certifying that two generated same-source,
same-target route records are distinguishable by a non-reducibility or
coherence-mismatch certificate. Take \(\mathrm{Level}_3\) on the
verification level, and let \(\pi_3(D_3)\) be the derived active
projection determined by
\((C_3,\Theta_3,W_3,\mathrm{Level}_3)\). The routine verifications
\(v_3^{\mathrm{fin}}, v_3^{\mathrm{cl}}, v_3^{\mathrm{aud}}\) then
follow. The nonclaim \(n_3\) excludes conflation with
\Pone{} route undefinedness and with direct cross-level comparison.
\end{proof}

\begin{proposition}[\Pfour{} lower bound]\label{prop:p4-witness}
There exists \(D_4 \in \FATCD{}\) with a derived active projection
\(\pi_4(D_4)\) witnessing \Pfour{}.\footnote{The Lean development establishes this for one concrete description \texttt{witnessDescription Role.P4} under parametric hypotheses (alignment, threshold attachment, and finite/closed/audited verifications of that fixed witness; the attached level is fixed to the verification level rather than the structural-categorical level of the prose); the narrowing is recorded in Appendix~\ref{app:mechanization} (\texttt{SixBirds.p4\_witness}).}
\end{proposition}

\begin{proof}
Let \(D_4\) carry index records \(i_0 < i_1 < i_2\) and transition
records \(t_{01}, t_{12}\) linking them, together with a refinement
record connecting two levels of indexing with explicit transition laws,
and closure records for the endpoints and refinement references used by
those transitions. Attach
\(\Theta_4 = \mathrm{Ann}_\Theta[\text{nontrivial stage/refinement}]\)
to the cluster \(C_4\) comprising these records, and record \(W_4\) as
the schema certifying a non-identity
order, transition, refinement, or stage tower. Take
\(\mathrm{Level}_4\) on the structural-categorical level, and let
\(\pi_4(D_4)\) be the derived active projection determined by
\((C_4,\Theta_4,W_4,\mathrm{Level}_4)\). The routine verifications
\(v_4^{\mathrm{fin}}, v_4^{\mathrm{cl}}, v_4^{\mathrm{aud}}\) then
follow. The nonclaim \(n_4\) excludes bare labels and arbitrary
relations or transitions without laws.
\end{proof}

\begin{proposition}[\Pfive{} lower bound]\label{prop:p5-witness}
There exists \(D_5 \in \FATCD{}\) with a derived active projection
\(\pi_5(D_5)\) witnessing \Pfive{}.\footnote{The Lean development establishes this for one concrete description \texttt{witnessDescription Role.P5} under parametric hypotheses (alignment over the same raw-kind cluster used for P1, threshold attachment, and finite/closed/audited verifications of that fixed witness); the narrowing is recorded in Appendix~\ref{app:mechanization} (\texttt{SixBirds.p5\_witness}).}
\end{proposition}

\begin{proof}
Let \(D_5\) carry distinct object records \(x \ne y\), a relation
record \(x \sim y\), a packaging map record
\(\mathrm{pack}\) with \(\mathrm{pack}(x) = \mathrm{pack}(y)\), and a
comparison record providing quotient-validity evidence. Attach
\(\Theta_5 = \mathrm{Ann}_\Theta[\text{nontrivial quotient packaging}]\)
to the cluster \(C_5\) comprising these records, and record \(W_5\) as
the schema certifying
that distinct raw objects are related and packaged together by a valid
packaging map with quotient-validity evidence. Take
\(\mathrm{Level}_5\) on the verification level, and let \(\pi_5(D_5)\)
be the derived active projection determined by
\((C_5,\Theta_5,W_5,\mathrm{Level}_5)\). The routine verifications
\(v_5^{\mathrm{fin}}, v_5^{\mathrm{cl}}, v_5^{\mathrm{aud}}\) then
follow. The nonclaim \(n_5\) excludes identity packaging and
weak-distinctness witnesses.
\end{proof}

\begin{proposition}[\Psix{} lower bound]\label{prop:p6-witness}
There exists \(D_6 \in \FATCD{}\) with a derived active projection
\(\pi_6(D_6)\) witnessing \Psix{}.\footnote{The Lean development establishes this for one concrete description \texttt{witnessDescription Role.P6} under parametric hypotheses (alignment over event/provenance/replay/check records, threshold attachment, and finite/closed/audited verifications of that fixed witness); the narrowing is recorded in Appendix~\ref{app:mechanization} (\texttt{SixBirds.p6\_witness}).}
\end{proposition}

\begin{proof}
Let \(D_6\) carry event records, a trace record, a provenance record
referencing the events, and a replay or check record. Record an audit
entry that references the event, tracks provenance and dependency, and
supports a replay or check relation. Attach
\(\Theta_6 = \mathrm{Ann}_\Theta[\text{nontrivial audit schema}]\) to
the cluster \(C_6\) comprising these records, and record \(W_6\) as the
schema certifying that
the audit record contains event-reference, provenance/dependency, and
replay/check structure. Take \(\mathrm{Level}_6\) on the verification
level, and let \(\pi_6(D_6)\) be the derived active projection
determined by \((C_6,\Theta_6,W_6,\mathrm{Level}_6)\). The routine
verifications \(v_6^{\mathrm{fin}}, v_6^{\mathrm{cl}},
v_6^{\mathrm{aud}}\) then follow. The nonclaim \(n_6\) excludes passive
logs, certificate-only content, and partial audit records.
\end{proof}

Propositions~\ref{prop:p1-witness}--\ref{prop:p6-witness} supply, for
each \(i \in \{1,\ldots,6\}\), an admissible witness tuple of the
form in Definition~\ref{def:common-witness-data}. This proves
Theorem~\ref{thm:six-role-lower-bounds}.

\begin{remark}[What the lower bounds do and do not establish]\label{rem:lower-bounds-scope}
Theorem~\ref{thm:six-role-lower-bounds} asserts that each of the six
roles is instantiated somewhere in \FATCD{}. It does not prove
exact-six, does not rule out candidate seventh-role threats, does not
prove decomposition, and does not prove uniqueness of any later
decomposition/exhaustion record. It also does not assert that any
single description activates all six channels simultaneously. Those
questions are handled in
Sections~\ref{sec:upper-bound-classification}
and~\ref{sec:decomposition-exhaustion}.
\end{remark}
 \section{Upper-bound classification: no seventh irreducible role}\label{sec:upper-bound-classification}




The lower bounds of Section~\ref{sec:lower-bounds} rule out the
trivialization in which some primitive role lacks an admissible
witness in \FATCD{}. The upper-bound direction addresses the
complementary question: whether some additional, seventh closure role
is needed in the covered scope. This section shows that no such role
is needed. The argument proceeds in three steps. First, we state the
alignment rules that constrain how raw features may project to
\Pone{}, \ldots, \Psix{}. Second, we close the three high-priority
candidate seventh-role threats: probability and uncertainty, residual
causality, and residual agency. Third, we record the upper-bound
classification theorem that collects these local classifications into
a single covered-scope statement.

\subsection{Role-projection alignment rules}\label{subsec:alignment-rules}

The role meanings of Definition~\ref{def:primitive-roles} carry
alignment rules constraining which raw features may be interpreted as
active projections of each role. These rules refine the corrected
alignments of Remark~\ref{rem:repaired-alignments} and form the
discipline under which the classification theorem is proved.

\begin{proposition}[Role-projection alignment]\label{prop:alignment-rules}
Under honest bookkeeping and the witness discipline of
Definition~\ref{def:common-witness-data}, each role projection is
constrained as follows.
\begin{itemize}[topsep=2pt,itemsep=1pt]
  \item \Pone{}: admissible only when the candidate cluster carries a
    selected update, a quotient or relation record, and a witness of
    descent compatibility, descent obstruction, or update-vs-quotient
    mismatch. Generic semantics, generic predicate meaning, generic
    selection criteria, and generic agency are not admissible as
    \Pone{} projections in the absence of such a witness.
  \item \Ptwo{}: admissible only when the candidate carries a macro
    specification, a realization relation, specification-satisfaction
    data, and evidence of representability or non-representability.
    Macro-level predicates or objectives count only insofar as they are
    treated as content to be represented. Mere quotient equality,
    arbitrary relations without representability structure, and
    empirical observations without level and threshold data are not
    admissible as \Ptwo{} projections.
  \item \Pthree{}: admissible only when the candidate carries enriched
    route grammar, generated same-source, same-target route terms, and
    route-comparison data with a non-reducibility or coherence-mismatch
    certificate. Generic transitions, generic causal arrows, generic
    path records, and route undefinedness arising from \Pone{}
    obstruction are not admissible as \Pthree{} projections unless the
    route/coherence threshold is explicitly present and the route
    comparison is specifically of the \Pthree{} type.
  \item \Pfour{}: admissible only when the candidate carries a stage,
    order, transition, refinement, staged-dynamics, or regime-ordering
    record with laws beyond bare labels. Bare labels, arbitrary
    relations without order, transition, or refinement law, quotient
    packaging, and audit bookkeeping are not admissible as \Pfour{}
    projections.
  \item \Pfive{}: admissible only when the candidate carries a quotient
    relation, a packaging map, the identification of distinct
    micro-elements, a same-package or quotient-object record, and
    quotient-validity evidence. Generic structural paths, derivations,
    topologies or geometries merely because they have paths, and causal
    chains are not admissible as \Pfive{} projections unless quotient
    or packaging data is present.
  \item \Psix{}: admissible only when the candidate carries an audit
    record with event reference, provenance or dependency trace, replay
    or check structure, and bookkeeping or checkability structure.
    Passive logs alone, arbitrary metadata, and certificate-only
    content without audit/provenance/replay/check structure are not
    admissible as \Psix{} projections.
\end{itemize}\footnote{Tracked in the formalization harness as a conditional schema (\texttt{SixBirds.alignment\_rules}); see Appendix~\ref{app:mechanization}.}
\end{proposition}

\begin{proof}
The constraints are not fresh stipulations; each is read off the role
meaning fixed in Definition~\ref{def:primitive-roles} together with the
witness discipline of Definition~\ref{def:common-witness-data}. By that
definition an admissible projection \(\pi_i(D)\) carries a witness schema
\(W_i\) that is, by requirement, the witness schema for the role meaning
of \(P_i\); a candidate cluster therefore qualifies as an active \(P_i\)
projection only if it supplies data of the type that the meaning of
\(P_i\) demands, and only then. We unpack this role by role.

For \Pone{}, the fixed meaning is update-vs-quotient compatibility and
descent obstruction. The witness type for that meaning is exactly a
selected update together with a quotient or relation record and a
certificate of descent compatibility, descent obstruction, or
update-vs-quotient mismatch; absent such a witness the candidate
exhibits no descent datum to compare, so it cannot be an active \Pone{}
projection. In particular generic semantics, generic predicate meaning,
generic selection criteria, and generic agency carry no descent witness
and are inadmissible, which is the corrected alignment of
Remark~\ref{rem:repaired-alignments}.

For \Ptwo{}, the fixed meaning is macro-specification representability.
Its witness type is a macro specification, a realization relation,
specification-satisfaction data, and a representability or
non-representability certificate; a macro-level predicate or objective
enters only when treated as content to be represented, since only then
does it instantiate the representability relation. Mere quotient
equality is \Pfive{} data, an arbitrary relation lacking
representability structure supplies no realization relation, and
empirical observations lacking level and threshold data fail the witness
discipline; none is an admissible \Ptwo{} projection.

For \Pthree{}, the fixed meaning is enriched route/coherence mismatch,
and \Pthree{} is not a directionality certificate. Its witness type is
enriched route grammar, generated same-source same-target route terms,
and route-comparison data carrying a non-reducibility or
coherence-mismatch certificate. A candidate without such a route
comparison has nothing of the \Pthree{} type to compare. Generic
transitions and generic causal arrows are \Pfour{} data, generic path
records are not route-comparison data, and route undefinedness arising
from a \Pone{} obstruction is a \Pone{} phenomenon; none projects to
\Pthree{} unless the route/coherence threshold is explicitly present and
the comparison is specifically of the \Pthree{} type. This is the
corrected alignment that generic transition is not \Pthree{}.

For \Pfour{}, the fixed meaning is stage, order, transition, and
refinement structure. Its witness type is a stage, order, transition,
refinement, staged-dynamics, or regime-ordering record carrying laws
beyond bare labels. Bare labels, arbitrary relations without an order,
transition, or refinement law, quotient packaging (which is \Pfive{}
data), and audit bookkeeping (which is \Psix{} data) supply no such law
and are inadmissible as \Pfour{} projections.

For \Pfive{}, the fixed meaning is quotient packaging. Its witness type
is a quotient relation, a packaging map, the identification of distinct
micro-elements, a same-package or quotient-object record, and
quotient-validity evidence. Generic structural paths, derivations,
topologies or geometries that merely carry paths, and causal chains
supply no quotient relation or packaging map, so they are inadmissible
absent quotient or packaging data; this is the corrected alignment that
generic path or derivation is not \Pfive{}.

For \Psix{}, the fixed meaning is audit, provenance, replay, and
checkability. Its witness type is an audit record with event reference,
a provenance or dependency trace, a replay or check structure, and
bookkeeping or checkability structure. Passive logs without check
structure, arbitrary metadata, and certificate-only content lacking
audit/provenance/replay/check structure supply no checkable audit datum
and are inadmissible as \Psix{} projections.

In each case the admissibility constraint is forced by the fixed role
meaning through the witness discipline: a projection is admissible only
with the witness type that the meaning of the role fixes.
\end{proof}

\begin{remark}[Three corrected alignments, held fixed]\label{rem:three-repaired-alignments}
Proposition~\ref{prop:alignment-rules} fixes three alignments that
are held throughout: generic transition data does not project to
\Pthree{}, generic path or derivation data does not project to
\Pfive{}, and generic semantic or interpretation data does not project
to \Pone{}. These are carried into every step of the classification
theorem below.
\end{remark}

\subsection{High-priority threat closures}\label{subsec:threat-closures}

Three classes of candidate seventh-role feature require particular
scrutiny: probability and uncertainty, residual causality in the
sense of the interventional and counterfactual
hierarchy~\citep{Pearl2009}, and residual agency in the planning
sense~\citep{Bratman1987}. Each closure is recorded as a proposition.
In every case the candidate is classified under the alignment rules
of Proposition~\ref{prop:alignment-rules} together with the residual
scheme of Definition~\ref{def:residual-scheme}; no candidate is
dismissed by definition. These closures remove the remaining
high-priority blockers before the general classification theorem.

\begin{proposition}[Closure of probability and uncertainty]\label{prop:probability-closure}
Every admissible probability or uncertainty feature of a description
\(D \in \FATCD{}\) whose content takes one of the recognized forms below
classifies, within the covered scope, as one of
\begin{itemize}[topsep=2pt,itemsep=1pt]
  \item a \Ptwo{} projection, when the feature is a macro-specification
    or representability condition;
  \item a \Pfour{} projection, when the feature is a stochastic
    transition or regime structure carrying valid stage or transition
    data;
  \item a \Psix{} projection, when the feature is an audited
    uncertainty-bookkeeping record;
  \item an activation-threshold refinement, when the uncertainty data
    modulates when an already admissible role projection activates
    without supplying a new witness schema;
  \item an empirical-bridge annotation, when the feature expresses
    measurement uncertainty bridged to a substrate observable;
  \item or a composite of the foregoing projections and annotations when
    the admissible bookkeeping carries several of these aspects at once.
\end{itemize}
Probability or uncertainty content of any of these recognized forms therefore
requires no seventh role; content of none of these forms is governed by the
recognition hypothesis (Definition~\ref{def:recognition-hypothesis}).\footnote{The Lean declaration is a trivial single-record harness check: \texttt{classifyResidual} maps \texttt{RawKind.probability} to \texttt{composite}, whose \texttt{CoveredResidualClass} branch is \texttt{True}; it does not establish the paper's universal probability/uncertainty closure. See Appendix~\ref{app:mechanization} (\texttt{SixBirds.probability\_closure}).}
\end{proposition}

\begin{proof}
Let \(r\) be an admissible probability or uncertainty feature of
\(D \in \FATCD{}\), and let it pass through the residual scheme of
Definition~\ref{def:residual-scheme} under the alignment rules of
Proposition~\ref{prop:alignment-rules}. By admissibility \(r\) carries
threshold and bookkeeping data; we case on the content that data fixes.

If the probabilistic content is a macro specification or a
representability or non-representability condition on a description,
then it supplies the realization-relation and representability witness
of \Ptwo{}, so by Proposition~\ref{prop:alignment-rules} it is a
\Ptwo{} projection. If instead the content is a stochastic transition,
regime, or staged structure carrying a valid stage, order, transition,
or refinement law, then it supplies the \Pfour{} witness and is a
\Pfour{} projection. If the content is an audited record of uncertainty
bookkeeping --- event reference, provenance, replay or check structure
--- it supplies the \Psix{} witness and is a \Psix{} projection. If the
uncertainty data does not supply any new witness schema but only
modulates when an already admissible projection activates, then by the
witness discipline it adds no role content and is an
activation-threshold refinement. If it expresses measurement
uncertainty bridged to a substrate observable, it carries an
empirical-bridge annotation rather than a role witness and is so
classified. Finally, when the admissible bookkeeping simultaneously
carries several of these aspects, \(r\) decomposes into the
corresponding projections and annotations and is a composite of the
foregoing. These are the recognized forms of probability and uncertainty content, and
each reduces to the six roles as stated; a probability feature whose content is
of none of these forms --- for instance a bare normalization or conservation
constraint --- is not classified here and falls under the recognition hypothesis
(Definition~\ref{def:recognition-hypothesis}).
A feature whose probabilistic content lies outside the covered scope is
classified as an outside-domain structure by
Definition~\ref{def:residual-scheme}, not absorbed by default. The
closure is therefore scoped, as stated.
\end{proof}

\begin{proposition}[Closure of residual causality]\label{prop:causality-closure}
Every admissible residual-causality feature of a description
\(D \in \FATCD{}\) whose content takes one of the recognized forms below
classifies, within the covered scope, as one of
\begin{itemize}[topsep=2pt,itemsep=1pt]
  \item a \Pfour{} projection, as the primary classification when the
    feature is a directed transition, refinement, or staged-dependence
    record;
  \item a \Pthree{} projection, only when the feature is specifically a
    route/coherence comparison between generated same-source,
    same-target routes;
  \item a \Psix{} projection, when the feature is a provenance or
    dependency audit;
  \item a composite of \Pfour{} and \Psix{}, when directed dependence
    and provenance audit occur together in the same admissible feature;
  \item an outside-domain structure, when the feature requires
    metaphysical causation not amenable to honest bookkeeping.
\end{itemize}
Classification as \Pfive{} is inadmissible in the absence of quotient
packaging. Residual-causality content of these recognized forms therefore requires no
seventh role; content of none of these forms is governed by the recognition
hypothesis (Definition~\ref{def:recognition-hypothesis}).\footnote{The Lean declaration is a trivial single-record harness check: \texttt{classifyResidual} maps \texttt{RawKind.causality} to \texttt{composite}, whose \texttt{CoveredResidualClass} branch is \texttt{True}; the \Pfive{}-inadmissibility content of the proposition is absent and no universal residual-causality closure is established. See Appendix~\ref{app:mechanization} (\texttt{SixBirds.causality\_closure}).}
\end{proposition}

\begin{proof}
Let \(r\) be an admissible residual-causality feature of
\(D \in \FATCD{}\), causality being taken in the interventional and
counterfactual sense of~\citet{Pearl2009}, and pass it through the
residual scheme of Definition~\ref{def:residual-scheme} under
Proposition~\ref{prop:alignment-rules}. A causal feature presents a
directed dependence, so we case on how that directedness is carried.

If \(r\) records a directed transition, a refinement, or a
staged-dependence law --- the typical form of admissible directed
dependence in \FATCD{} --- it supplies the order/transition/refinement
witness of \Pfour{}, so by Proposition~\ref{prop:alignment-rules} it is
a \Pfour{} projection; this is the primary classification. The single
exception in which directedness routes elsewhere is when \(r\) is
specifically a route/coherence comparison between generated
same-source, same-target routes carrying a non-reducibility or
coherence-mismatch certificate; only then does it supply the \Pthree{}
witness and classify as a \Pthree{} projection. We must not assign a
generic causal arrow to \Pthree{}, since by the alignment rules a
generic transition is not a route comparison of the \Pthree{} type. If
\(r\) is instead a provenance or dependency audit --- event reference
with a provenance or dependency trace --- it supplies the \Psix{}
witness and is a \Psix{} projection. When directed dependence and
provenance audit occur together in the same admissible feature, \(r\)
decomposes into both and is a composite of \Pfour{} and \Psix{}.
Finally, if \(r\) asserts metaphysical causation not amenable to honest
bookkeeping --- carrying no finite directed-dependence, route, or audit
witness --- it lies outside the covered scope and is classified as an
outside-domain structure.

Classification as a \Pfive{} projection is not available here: by
Proposition~\ref{prop:alignment-rules} a \Pfive{} projection requires a
quotient relation and packaging map, and a causal feature carrying no
quotient packaging supplies none; the directedness of a causal chain is
not packaging data. Hence the \Pfive{}-inadmissibility caveat holds.
The cases above are the recognized forms of residual-causality content; a
causality feature whose content is of none of these forms is not classified here
and falls under the recognition hypothesis
(Definition~\ref{def:recognition-hypothesis}). The closure is scoped and makes no
scope-free claim.
\end{proof}

\begin{proposition}[Closure of residual agency]\label{prop:agency-closure}
Every admissible residual-agency feature of a description
\(D \in \FATCD{}\) whose content takes one of the recognized forms below
classifies, within the covered scope, as one of
\begin{itemize}[topsep=2pt,itemsep=1pt]
  \item a composite of \Ptwo{}, \Pfour{}, and \Psix{}, when the feature
    decomposes into an objective or specification to be represented, a
    transition, regime, or staged selection event, and an audited
    decision-provenance record;
  \item an empirical-bridge annotation, when observed agency serves as
    substrate evidence for a bridged claim;
  \item an outside-domain structure, when irreducible intentional or
    teleological content is asserted beyond what admissible
    bookkeeping reduces to \Ptwo{}, \Pfour{}, and \Psix{}.
\end{itemize}
Classification as \Pone{} is inadmissible in the absence of an explicit
update-vs-quotient descent obstruction. Residual-agency content of these recognized forms therefore requires no
seventh role; content of none of these forms is governed by the recognition
hypothesis (Definition~\ref{def:recognition-hypothesis}).\footnote{The Lean declaration is a trivial single-record harness check: \texttt{classifyResidual} maps \texttt{RawKind.agency} to \texttt{empiricalBridgeAnnotation}, whose \texttt{CoveredResidualClass} branch is \texttt{True}; the \Pone{}-inadmissibility content is absent and no universal residual-agency closure is established. See Appendix~\ref{app:mechanization} (\texttt{SixBirds.agency\_closure}).}
\end{proposition}

\begin{proof}
Let \(r\) be an admissible residual-agency feature of
\(D \in \FATCD{}\), agency being taken in the planning sense
of~\citet{Bratman1987}, and pass it through the residual scheme of
Definition~\ref{def:residual-scheme} under
Proposition~\ref{prop:alignment-rules}. An agency feature in admissible
form presents an objective pursued through selection and recorded for
later inspection; we case on whether that content reduces to role
witnesses or asserts more.

If \(r\) decomposes into an objective or specification to be
represented, a transition, regime, or staged selection event, and an
audited decision-provenance record, then each component supplies the
corresponding witness: the objective-as-content supplies the \Ptwo{}
representability witness, the selection event supplies the \Pfour{}
transition or staged-dynamics witness, and the decision-provenance
record supplies the \Psix{} audit witness. By
Proposition~\ref{prop:alignment-rules} \(r\) is then a composite of
\Ptwo{}, \Pfour{}, and \Psix{}. If instead the observed agency serves
only as substrate evidence for a bridged claim, it carries an
empirical-bridge annotation and is so classified. If \(r\) asserts
irreducible intentional or teleological content beyond what admissible
bookkeeping reduces to \Ptwo{}, \Pfour{}, and \Psix{}, then it carries
no finite role witness and lies outside the covered scope, classified as
an outside-domain structure.

Classification as a \Pone{} projection is not available here merely
because agency selects an update: by
Proposition~\ref{prop:alignment-rules} a \Pone{} projection requires a
selected update together with a quotient or relation record and an
explicit descent-compatibility, descent-obstruction, or
update-vs-quotient mismatch witness. A selection event carrying no such
descent obstruction supplies no \Pone{} witness, so its update content
classifies under \Pfour{} (and, where audited, \Psix{}) rather than
\Pone{}. This is the stated \Pone{}-inadmissibility caveat. The cases
above are the recognized forms of residual-agency content; content of none of
these forms falls under the recognition hypothesis
(Definition~\ref{def:recognition-hypothesis}). The closure is scoped and asserts
nothing beyond it.
\end{proof}

\begin{remark}[Scope of the closures]\label{rem:threat-closures-scope}
Propositions~\ref{prop:probability-closure}--\ref{prop:agency-closure}
close the three threats within the covered scope; they do not make
scope-free claims. A feature that falls outside the covered scope is
classified as an outside-domain structure via
Definition~\ref{def:residual-scheme}, not absorbed into a primitive role
by default.
\end{remark}

\subsection{The recognition hypothesis and the classification theorem}\label{subsec:classification-theorem}

\begin{definition}[Recognition hypothesis]\label{def:recognition-hypothesis}
A description \(D \in \FATCD{}\) satisfies the \emph{recognition hypothesis}
when every residual feature \(r \in \Residual{D}\) is \emph{recognized}: the
closure content of \(r\) either supplies the witness type of one of the six
role meanings of Definition~\ref{def:primitive-roles} --- so that \(r\), or a
finite composite of its constituents, projects to \Pone{}, \ldots, \Psix{}
under Proposition~\ref{prop:alignment-rules} --- or instantiates one of the six
annotation kinds of Definition~\ref{def:annotations}, or is substrate already
accounted for by the audit discipline of \FATCD{}. Equivalently, \(D\) carries
no finite audited residual feature whose closure content is irreducible to the
six role meanings, and in particular none whose role-theoretic content emerges
only in an assembled cluster and is not the join of its constituents' role
labels.
\end{definition}

\begin{remark}[The recognition hypothesis is assumed, not derived]\label{rem:recognition-open}
The recognition hypothesis is an explicit assumption on \(D\), not a consequence
of the \FATCD{} definitions. Two kinds of finite, audited, in-domain feature are
not known to be recognized: a feature carrying a conservation or balance law over
declared records --- a finite resource invariant of the kind admitted by
Definition~\ref{def:finite} --- whose content supplies none of the six witness
types, and a feature whose closure content is an emergent obstruction or holonomy
of an assembled cluster, not reducible to the role labels of its constituents.
Whether every finite audited feature is recognized --- equivalently, whether the
six roles are \emph{recognition-complete} for \FATCD{} --- is left open and is
recorded with the endogenous-instrument meta-limit of
Sections~\ref{subsec:meta-limit} and~\ref{subsec:future-meta-limit}. The three
threat closures of
Propositions~\ref{prop:probability-closure}--\ref{prop:agency-closure} establish
the recognition hypothesis unconditionally for the three highest-priority threat
classes; the classification theorem below assumes it in general.
\end{remark}

The alignment rules of Proposition~\ref{prop:alignment-rules}, the
threat closures of
Propositions~\ref{prop:probability-closure}--\ref{prop:agency-closure},
and the residual scheme of Definition~\ref{def:residual-scheme}
combine to yield the upper-bound classification theorem: alignment
rules fix the admissible role meanings, threat closures rule out the
three misclassifications that otherwise block the argument, and the
theorem, under the recognition hypothesis
(Definition~\ref{def:recognition-hypothesis}), records the classification
codomain for every residual feature.

\begin{theorem}[Upper-bound classification]\label{thm:upper-bound-classification}
Let \(D \in \FATCD{}\) satisfy the recognition hypothesis
(Definition~\ref{def:recognition-hypothesis}), and let \(r\) be a residual
feature of \(D\) in the sense of Definition~\ref{def:residual-scheme}. Then
\(r\) classifies as one of
\begin{enumerate}[label=(\roman*),topsep=2pt,itemsep=1pt]
  \item already absorbed into an admissible \Pone{}, \Ptwo{}, \Pthree{},
    \Pfour{}, \Pfive{}, or \Psix{} projection under the alignment rules of
    Proposition~\ref{prop:alignment-rules};
  \item a composite of such projections;
  \item a presentation variant of an already classified feature;
  \item a level shift of an already classified feature;
  \item an activation-threshold refinement;
  \item an instrument-visibility artifact;
  \item an empirical-bridge annotation;
  \item an outside-domain structure;
  \item already-covered bookkeeping or annotation data.
\end{enumerate}
Equivalently, within the covered scope, categories (x) and (xi) of
Definition~\ref{def:residual-scheme}, namely the unresolved-threat and
genuine-seventh-role categories, are empty: no covered candidate
seventh-role feature remains as an unresolved threat, and no covered
candidate seventh-role feature requires a genuine seventh irreducible
typed closure role.\footnote{The Lean declaration is a classifier-by-construction check: every \texttt{RawKind} branch of \texttt{classifyResidual} lands in a class whose \texttt{CoveredResidualClass} branch is \texttt{True}, and the uncovered constructors \texttt{unresolvedThreat}/\texttt{genuineSeventhRole} are never produced; ``no seventh role'' holds by construction rather than by derivation from feature content. See Appendix~\ref{app:mechanization} (\texttt{SixBirds.upper\_bound\_classification}).}
\end{theorem}

\begin{proof}
Let \(r\) be a residual feature of \(D \in \FATCD{}\) in the covered
scope, in the sense of Definition~\ref{def:residual-scheme}: a raw
record or annotation in \(\Raw{D}\) whose role-theoretic disposition is
not fixed by its raw-grammar kind alone. By
Definition~\ref{def:residual-scheme} the scheme assigns \(r\) a single
label \(\operatorname{Classify}(r)\) drawn from (i)--(xi); we show that,
within the covered scope, applying the alignment discipline of
Proposition~\ref{prop:alignment-rules} together with the three threat
closures forces that label into (i)--(ix), leaving (x) and (xi) empty.

\emph{Disposition under the alignment discipline.} Since \(r\) is
admissible, it carries the threshold, witness, level, and bookkeeping
data of Definition~\ref{def:common-witness-data}. We test that data
against the six fixed role meanings using the alignment rules of
Proposition~\ref{prop:alignment-rules}, which (by that proposition) admit
a projection to \(P_i\) exactly when \(r\) supplies the witness type that
the meaning of \(P_i\) fixes. There are two possibilities. Either the
witness, threshold, and level data of \(r\) align with one of the fixed
role meanings, in which case \(r\) is absorbed into an admissible
\Pone{}, \Ptwo{}, \Pthree{}, \Pfour{}, \Pfive{}, or \Psix{} projection
and \(\operatorname{Classify}(r) = \) (i); or the data of \(r\) require
several such aligned components at once --- for instance an objective to
be represented carried over a staged selection with an audit trail ---
in which case \(r\) decomposes into those projections and
\(\operatorname{Classify}(r) = \) (ii). The corrected alignments of
Remark~\ref{rem:repaired-alignments}, carried in through
Remark~\ref{rem:three-repaired-alignments}, are used here to keep these
absorptions honest: a generic transition does not slip into (i) as a
\Pthree{} projection, a generic path or derivation does not slip in as a
\Pfive{} projection, and generic semantics does not slip in as a \Pone{}
projection.

\emph{The bookkeeping and presentation categories.} If \(r\) is not a
fresh role witness but a re-presentation of an already classified
feature, it is a presentation variant, \(\operatorname{Classify}(r) = \)
(iii); if it restates an already classified feature at a different level
of the trichotomy of Definition~\ref{def:level-trichotomy}, it is a
level shift, (iv). If its uncertainty or threshold data only modulates
when an already admissible projection activates, without supplying a new
witness schema, it is an activation-threshold refinement, (v). If it is
visible or suppressed only relative to a declared instrument, it is an
instrument-visibility artifact, (vi). If it bridges a substrate
observable to an admissible claim, it is an empirical-bridge annotation,
(vii). If it requires non-finite or unbookkept semantics not amenable to
honest bookkeeping, it falls outside the covered scope and is an
outside-domain structure, (viii). If it is bookkeeping or annotation data
already accounted for by the audit discipline of \FATCD{}, it is
already-covered bookkeeping or annotation data, (ix). These are the dispositions available to a recognized feature.

\emph{The three high-priority threats.} The classes most apt to demand a
seventh role are recognized unconditionally. A probability or uncertainty
feature of one of the recognized forms is, by
Proposition~\ref{prop:probability-closure}, a \Ptwo{}, \Pfour{}, or \Psix{}
projection, an activation-threshold refinement, an empirical-bridge annotation,
or a composite of these --- all in (i)--(ix). A residual-causality feature of a
recognized form is, by Proposition~\ref{prop:causality-closure}, a \Pfour{},
\Pthree{}, or \Psix{} projection, a \Pfour{}/\Psix{} composite, or an
outside-domain structure, with \Pfive{} excluded for want of quotient packaging.
A residual-agency feature of a recognized form is, by
Proposition~\ref{prop:agency-closure}, a \Ptwo{}/\Pfour{}/\Psix{} composite, an
empirical-bridge annotation, or an outside-domain structure, with \Pone{}
excluded for want of a descent obstruction. The three closures thus establish the
recognition hypothesis for the three highest-priority threat classes without
assuming it.

\emph{Conclusion.} Under the recognition hypothesis
(Definition~\ref{def:recognition-hypothesis}), every residual feature \(r\) of
\(D\) is recognized and hence receives a label in (i)--(ix): its content supplies
a role witness (cases (i)--(ii)), instantiates an annotation kind (cases
(iii)--(vii)), or is already-covered substrate (case (ix)), and only non-finite
or unbookkept content is sent to the outside-domain category~(viii). Since every residual feature receives a label in (i)--(ix) by the foregoing,
category~(x), the unresolved-threat category, is empty; and since the recognition hypothesis
assumes that no finite audited in-domain feature has content irreducible to the
six roles, category~(xi), the genuine-seventh-role category, is empty. The
emptiness of~(xi) is the content of the recognition hypothesis, not a consequence
of the \FATCD{} definitions: by Remark~\ref{rem:recognition-open} a finite
audited conservation law or an emergent composite obstruction would
populate~(xi), and whether any such feature exists is the open
recognition-completeness question. The theorem holds exactly to the strength of
its hypothesis.
\end{proof}

\begin{remark}[Scope of the theorem]\label{rem:classification-scope}
Theorem~\ref{thm:upper-bound-classification} is conditional on the recognition
hypothesis (Definition~\ref{def:recognition-hypothesis}) and scoped to \FATCD{}.
It does not derive that no finite audited feature requires a seventh role; it
assumes this, and discharges the assumption unconditionally only for the three
threat classes of
Propositions~\ref{prop:probability-closure}--\ref{prop:agency-closure}. It does
not claim that no seventh role can arise outside \FATCD{}, that the alignment
rules are exhaustive for richer domains, or that the recognition hypothesis holds
in general. Nor does it prove exact-six on its own: the decomposition/exhaustion
theorem of Section~\ref{sec:decomposition-exhaustion} remains necessary. The open
recognition-completeness question and broader questions are deferred to
Sections~\ref{subsec:meta-limit} and~\ref{sec:future-work}.
\end{remark}
 \section{Decomposition and exhaustion}\label{sec:decomposition-exhaustion}




The lower-bound theorem of Section~\ref{sec:lower-bounds} and the
upper-bound classification of
Section~\ref{sec:upper-bound-classification} reduce the scoped
exact-six question to a decomposition-existence problem: given
\(D \in \FATCD{}\), does \(D\) admit a well-formed record combining
the six role channels, their statuses, and the classification of any
raw features not absorbed into active projections? This section
answers affirmatively and records the channel-versus-activation
discipline that the answer requires. The lower bounds show that each
of the six channels is materially realizable in the domain; the
upper-bound classification shows that, under the recognition hypothesis, no
covered seventh role channel is needed.

\subsection{Decomposition records and channel statuses}\label{subsec:dec-records}

\begin{definition}[Channel status]\label{def:channel-status}
For each \(i \in \{1,\ldots,6\}\), the \emph{role channel}
\(\Chan{i}{D}\) of a description \(D \in \FATCD{}\) is the
\(P_i\)-labeled slot in the six-channel interface associated to \(D\),
regardless of whether that slot is active. A \emph{channel status} for
\(\Chan{i}{D}\) is exactly one of the five values
\begin{itemize}[topsep=2pt,itemsep=1pt]
  \item \(\activeproj\) --- an active derived role projection
    \(\pi_i(D)\) in the sense of
    Definition~\ref{def:role-projection} exists, with the full
    attachment data of
    Definition~\ref{def:threshold-attachment};
  \item \(\belowthr\) --- candidate raw features are present but the
    threshold data required by
    Definition~\ref{def:threshold-attachment} is absent;
  \item \(\collapsedbc\) --- the candidate cluster is a boundary case
    that collapses under the witness schema \(W_i\);
  \item \(\absentrec\) --- no relevant raw features for role \(P_i\) are
    present in \(\Raw{D}\), and the absence is itself recorded;
  \item \(\outsidescope\) --- the role-\(P_i\) content of \(D\) lies
    outside the covered scope.
\end{itemize}
We write \(\Status{i}{D}\) for the channel status of \(\Chan{i}{D}\).
\end{definition}

\begin{definition}[Decomposition/exhaustion record]\label{def:decomposition-record}
A \emph{decomposition/exhaustion record} for \(D \in \FATCD{}\) is a tuple
\[
\Dec{D} = \bigl(\Raw{D},\; (\Chan{i}{D}, \Status{i}{D})_{i=1}^{6},\;
\Residual{D},\; \mathrm{Classify},\; \mathcal{L},\; \mathcal{V},\;
\Audit{D},\; \mathcal{N}\bigr)
\]
where
\begin{itemize}[topsep=2pt,itemsep=1pt]
  \item the six role channels together with their statuses form the
    six-channel status sequence. If \(\Status{i}{D} = \activeproj\), the channel
    carries an active projection \(\pi_i(D)\) with its full attachment
    data. If \(\Status{i}{D}\) is inactive, the channel carries the
    corresponding below-threshold, collapsed-boundary,
    absent-with-record, or outside-scope status record;
  \item \(\Residual{D} \subseteq \Raw{D}\) is the residual feature set
    (the raw records and annotations not already accounted for by an
    active projection or an inactive status record), and
    \(\mathrm{Classify}\) assigns to each \(r \in \Residual{D}\) a label
    in the scheme of Definition~\ref{def:residual-scheme};
  \item \(\mathcal{L}\) collects the loss and lift records for any level
    crossings involved in projections or classifications;
  \item \(\mathcal{V}\) collects instrument-visibility and
    empirical-bridge annotations where these are asserted;
  \item \(\Audit{D}\) is the honest bookkeeping subrecord;
  \item \(\mathcal{N}\) is the explicit nonclaim record.
\end{itemize}
A decomposition/exhaustion record is \emph{well-formed} when each
\(\Status{i}{D}\) is drawn from the five values of
Definition~\ref{def:channel-status}, each active
\(\pi_i(D)\) is admissible in the sense of
Definitions~\ref{def:role-projection} and~\ref{def:threshold-attachment},
each inactive status record is honestly supported by the corresponding
raw-feature situation, and every raw feature is accounted for at the
channel/exhaustion level either by an active projection, by a status
record for its channel, or by a residual classification. The records in
\(\mathcal{L}\) and \(\mathcal{V}\) are auxiliary annotations on this
coverage, not additional role channels.
\end{definition}


\subsection{Existence, six-channel exactness, and residual exhaustion}\label{subsec:dec-theorems}

\begin{theorem}[Decomposition existence]\label{thm:decomposition-existence}
Every description \(D \in \FATCD{}\) admits a well-formed decomposition/exhaustion
record \(\Dec{D}\) in the sense of
Definition~\ref{def:decomposition-record}.\footnote{Tracked in the formalization harness as a conditional schema (\texttt{SixBirds.decomposition\_existence}); see Appendix~\ref{app:mechanization}.}
\end{theorem}

\begin{proof}
For each \(i \in \{1,\ldots,6\}\), inspect \(\Raw{D}\) for a cluster
aligned with the meaning of \(P_i\) under
Proposition~\ref{prop:alignment-rules}. If such a cluster carries the
threshold, witness, level, and audit data of
Definition~\ref{def:threshold-attachment}, record
\(\Status{i}{D} = \activeproj\) and an active projection
\(\pi_i(D)\) in the sense of Definition~\ref{def:role-projection};
otherwise record one of the four inactive statuses, determined by
which attachment data is missing, whether the cluster is a
boundary-case collapse, whether no relevant cluster exists, or
whether the role lies outside the covered scope. Six status entries
result, one per role. The raw features absorbed by these channel
records are removed from the residual stage; the remainder pass to
\(\Residual{D}\). By Definition~\ref{def:residual-scheme} the scheme assigns
each residual feature a single label drawn from (i)--(xi), so
\(\mathrm{Classify}\) is total on \(\Residual{D}\) and the record is well-formed
irrespective of which labels occur; under the recognition hypothesis
(Definition~\ref{def:recognition-hypothesis}),
Theorem~\ref{thm:upper-bound-classification} further places every label in
(ii)--(ix), with (x) and (xi) empty. The records \(\mathcal{L}\) and \(\mathcal{V}\) are
read from \(\Raw{D}\) and \(\Audit{D}\); the bookkeeping subrecord is
\(\Audit{D}\); and \(\mathcal{N}\) records the nonclaims carried
forward from the witness discipline. The resulting tuple is a
decomposition/exhaustion record, well-formed by construction.
Non-vacuity of the six channels follows from
Theorem~\ref{thm:six-role-lower-bounds}, which supplies, for each
channel, a description in which that channel is active.
\end{proof}

\begin{theorem}[Six-channel exactness and residual exhaustion]\label{thm:six-channel-exactness}
For every \(D \in \FATCD{}\) and every well-formed
\(\Dec{D}\) in the sense of
Definition~\ref{def:decomposition-record}:
\begin{enumerate}[label=(\roman*),topsep=2pt,itemsep=1pt]
  \item the channel sequence \((\Chan{i}{D})_{i=1}^{6}\) has length
    exactly six, one per primitive role \Pone{}, \ldots, \Psix{};
  \item every raw feature of \(D\) is accounted for at the
    channel/exhaustion level by an active projection \(\pi_i(D)\), a
    channel-status record for an inactive channel, or a residual
    classification under the scheme of
    Definition~\ref{def:residual-scheme}; the records in \(\mathcal{L}\)
    and \(\mathcal{V}\) are auxiliary annotations on this coverage;
  \item under the recognition hypothesis
    (Definition~\ref{def:recognition-hypothesis}) and within the covered scope,
    no raw feature is classified as an unresolved seventh-role threat or as a
    genuine seventh irreducible typed closure role
    (Theorem~\ref{thm:upper-bound-classification}); the three principal threat
    classes are closed unconditionally.
\end{enumerate}\footnote{Tracked in the formalization harness as a conditional schema (\texttt{SixBirds.six\_channel\_exactness}); see Appendix~\ref{app:mechanization}.}
\end{theorem}

\begin{proof}
Exactness of the channel sequence is built into the definition of
\(\Dec{D}\); there are six role channels by
Definition~\ref{def:decomposition-record}. Coverage of raw features is
the well-formedness condition of
Definition~\ref{def:decomposition-record} together with totality of
\(\mathrm{Classify}\) on \(\Residual{D}\), which holds by
Definition~\ref{def:residual-scheme}. Emptiness of categories (x) and (xi),
under the recognition hypothesis (Definition~\ref{def:recognition-hypothesis})
and within the covered scope, is
Theorem~\ref{thm:upper-bound-classification} directly.
\end{proof}

\subsection{Channel-versus-activation discipline}\label{subsec:channel-vs-activation}

Theorem~\ref{thm:six-channel-exactness} asserts six channels in every
decomposition/exhaustion record. It does not assert six active
projections in every record, nor does it assert that a well-formed
record is unique. The distinction between channel presence and
channel activation prevents six-channel exactness from being read as
an all-six-active claim.

\begin{proposition}[Channel-versus-activation discipline]\label{prop:channel-vs-activation}
Let \(D \in \FATCD{}\) and let \(\Dec{D}\) be a well-formed decomposition/exhaustion
record. The following are distinct statements:
\begin{enumerate}[label=(\roman*),topsep=2pt,itemsep=1pt]
  \item \emph{Six channels}: the channel sequence
    \((\Chan{i}{D})_{i=1}^{6}\) has length six with the fixed role names
    \Pone{}, \ldots, \Psix{}.
  \item \emph{Six active projections}: for each \(i\),
    \(\Status{i}{D} = \activeproj\) and the channel carries an active
    \(\pi_i(D)\).
\end{enumerate}
Theorem~\ref{thm:six-channel-exactness} establishes (i) for every
\(D \in \FATCD{}\). It does not establish (ii). In particular, an
admissible \(D \in \FATCD{}\) may carry any number of active projections
between zero and six; inactive channels are recorded with status
\(\belowthr\), \(\collapsedbc\), \(\absentrec\), or \(\outsidescope\)
and are not active witnesses for the corresponding roles.\footnote{Tracked in the formalization harness as a conditional schema (\texttt{SixBirds.channel\_vs\_activation}); see Appendix~\ref{app:mechanization}.}
\end{proposition}

\begin{proposition}[Decomposition records are not unique]\label{prop:non-uniqueness}
A description \(D \in \FATCD{}\) may admit more than one well-formed
decomposition/exhaustion record in the sense of
Definition~\ref{def:decomposition-record}. Distinct well-formed records
may differ in level assignments, choice of witness schema, residual
classification labels for features admitting more than one admissible
category under the alignment rules, or placement of loss/lift and
visibility records. Theorems~\ref{thm:decomposition-existence}
and~\ref{thm:six-channel-exactness} assert existence, not uniqueness.\footnote{The Lean development exhibits two concrete decomposition records of a single description that differ only in their \texttt{P1} channel status; the Lean statement does not assert that either record is well-formed, and the further sources of non-uniqueness in the prose (level assignments, witness-schema choice, residual labels, loss/lift placement) are not mechanized. The narrowing is recorded in Appendix~\ref{app:mechanization} (\texttt{SixBirds.non\_uniqueness}).}
\end{proposition}

\begin{remark}[What the decomposition theorems do and do not prove]\label{rem:dec-scope}
Theorems~\ref{thm:decomposition-existence}
and~\ref{thm:six-channel-exactness} establish the structural side of
the scoped exact-six claim: every \(D \in \FATCD{}\) admits a
six-channel decomposition/exhaustion record, with residual exhaustion in the
covered scope under the recognition hypothesis
(Definition~\ref{def:recognition-hypothesis}). Together with the lower-bound theorem
(Section~\ref{sec:lower-bounds}) and the upper-bound classification
(Section~\ref{sec:upper-bound-classification}), they prepare the
scoped exact-six theorem but do not replace the integrated-stress
check that follows. They do not assert that every description
activates all six roles, that the decomposition is canonical or
unique, or that the scoped exact-six program extends beyond \FATCD{};
those remain nonclaims.
\end{remark}

 \section{Integrated stress}\label{sec:integrated-stress}




With the theorem chain of Sections~\ref{sec:lower-bounds},
\ref{sec:upper-bound-classification}, and
\ref{sec:decomposition-exhaustion} in place, we record six stress
families that act as consistency checks against overclaims. These are
not proofs of the scoped exact-six target; they verify that the
theorem chain does not silently admit overclaims. Each family is a
guardrail that restates a nonclaim already secured by an earlier
result, so each is recorded as a remark. Their joint satisfaction is
the last requirement before the terminal statement in
Section~\ref{sec:scoped-exact-six}.

\subsection{Dependency stress}\label{subsec:dependency-stress}

\begin{remark}[Dependency stress]\label{prop:dependency-stress}
The scoped exact-six target is dependency-coherent only when the theorem
chain is assembled in the order
\[
\text{\FATCD{} domain} \;\rightarrow\; \text{lower bounds} \;\rightarrow\;
\text{upper-bound classification} \;\rightarrow\;
\text{decomposition/exhaustion}.
\]
Accordingly, the following are inadmissible:
\begin{itemize}[topsep=2pt,itemsep=1pt]
  \item asserting the target without the \FATCD{} domain of
    Section~\ref{sec:fatcd-domain};
  \item asserting the target without the lower-bound theorem of
    Section~\ref{sec:lower-bounds};
  \item asserting the target without the upper-bound classification
    theorem of Section~\ref{sec:upper-bound-classification};
  \item asserting the target without the decomposition and exhaustion
    theorems of Section~\ref{sec:decomposition-exhaustion};
  \item promoting the assembled target before the integrated stress
    check of Proposition~\ref{thm:integrated-stress}.
\end{itemize}
\end{remark}

\subsection{Circularity stress}\label{subsec:circularity-stress}

\begin{remark}[Circularity stress]\label{prop:circularity-stress}
Three circularity patterns are inadmissible under the theorem chain:
\begin{itemize}[topsep=2pt,itemsep=1pt]
  \item redefining \FATCD{} so that it already contains exactly six role
    slots by construction;
  \item declaring the absence of a seventh role by definition rather
    than by classification under Definition~\ref{def:residual-scheme};
  \item justifying decomposition existence by assuming scoped
    exact-six.
\end{itemize}
The non-circularity schema of Proposition~\ref{prop:non-circularity} and
the constructive proof of Theorem~\ref{thm:decomposition-existence}
block each pattern.
\end{remark}

\subsection{Channel-versus-activation stress}\label{subsec:channel-activation-stress}

\begin{remark}[Channel-versus-activation stress]\label{prop:channel-activation-stress}
The following conflations are inadmissible:
\begin{itemize}[topsep=2pt,itemsep=1pt]
  \item requiring that every \(D \in \FATCD{}\) activates all six role
    channels;
  \item treating an inactive channel of status \(\belowthr\) as an
    active witness for its role;
  \item treating an inactive channel of status \(\collapsedbc\) as an
    active witness;
  \item treating an inactive channel of status \(\absentrec\) as an
    active witness.
\end{itemize}
Proposition~\ref{prop:channel-vs-activation} enforces each distinction.
The same inactive-status guardrail also applies to
\(\outsidescope\).
\end{remark}

\subsection{Role-alignment stress}\label{subsec:role-alignment-stress}

\begin{remark}[Role-alignment stress]\label{prop:role-alignment-stress}
The three forbidden assignments of
Remark~\ref{rem:three-repaired-alignments} are inadmissible: generic
transition data does not project to \Pthree{}, generic path or
derivation data does not project to \Pfive{}, and generic semantic or
interpretation data does not project to \Pone{}. The alignment rules
of Proposition~\ref{prop:alignment-rules} enforce each prohibition.
\end{remark}

\subsection{Upper-bound residual stress}\label{subsec:upper-bound-residual-stress}

\begin{remark}[Upper-bound residual stress]\label{prop:residual-stress}
Under Theorem~\ref{thm:upper-bound-classification} and within the covered
scope, the following regressions are inadmissible:
\begin{itemize}[topsep=2pt,itemsep=1pt]
  \item probability or uncertainty features reopening as unresolved
    threats;
  \item residual-causality features reopening as unresolved threats;
  \item residual-agency features reopening as unresolved threats;
  \item any covered residual feature remaining unclassified under
    Definition~\ref{def:residual-scheme};
  \item any covered residual feature requiring a genuine seventh
    irreducible typed closure role.
\end{itemize}
Propositions~\ref{prop:probability-closure}--\ref{prop:agency-closure}
settle the three threat closures, and
Theorem~\ref{thm:upper-bound-classification}, under the recognition
hypothesis, delivers the emptiness of categories (x) and (xi).
\end{remark}

\subsection{Overclaim stress}\label{subsec:overclaim-stress}

\begin{remark}[Overclaim stress]\label{prop:overclaim-stress}
The following restatements of the scoped exact-six target are
inadmissible:
\begin{itemize}[topsep=2pt,itemsep=1pt]
  \item restatement as a full six-primitives algebra;
  \item restatement as a global primitive-independence theorem;
  \item restatement as unique or canonical decomposition;
  \item restatement as empirical realization;
  \item restatement as a claim about all emergence phenomena, or any
    domain broader than \FATCD{}.
\end{itemize}
Remark~\ref{rem:atlas-global-nonclaims},
Proposition~\ref{prop:nonclaim-closure}, and
Proposition~\ref{prop:non-uniqueness} carry the corresponding nonclaims
forward through the theorem chain. In particular, nonclaim closure
survives the exact-six assembly.
\end{remark}

\subsection{The integrated stress certification}\label{subsec:integrated-stress-theorem}

\begin{proposition}[Integrated stress certification]\label{thm:integrated-stress}
The theorem chain of Sections~\ref{sec:fatcd-domain},
\ref{sec:lower-bounds}, \ref{sec:upper-bound-classification}, and
\ref{sec:decomposition-exhaustion} satisfies each of the six stress
families recorded in the stress remarks of this section.
Dependency, circularity, channel-versus-activation, role-alignment,
upper-bound residual, and overclaim regressions are all blocked, and
the assembled target is ready for the terminal statement of
Section~\ref{sec:scoped-exact-six}.\footnote{Tracked in the formalization harness as a conditional schema (\texttt{SixBirds.integrated\_stress}) bundling the six stress conjuncts; see Appendix~\ref{app:mechanization}.}
\end{proposition}

\begin{proof}
We verify that each of the six guardrails is discharged by the earlier
result it rests on, so that none of the listed regressions can arise in
the assembled chain.

\emph{Dependency.} The components named in
Remark~\ref{prop:dependency-stress} are present and ordered: the
\FATCD{} domain is fixed in Section~\ref{sec:fatcd-domain}, the
lower-bound theorem in Section~\ref{sec:lower-bounds}
(Theorem~\ref{thm:six-role-lower-bounds}), the upper-bound
classification in Section~\ref{sec:upper-bound-classification}
(Theorem~\ref{thm:upper-bound-classification}), and the
decomposition and exhaustion theorems in
Section~\ref{sec:decomposition-exhaustion}
(Theorems~\ref{thm:decomposition-existence}
and~\ref{thm:six-channel-exactness}). Each later statement is stated
over the output of the previous one, so the target cannot be asserted
while skipping any link, and no link circularly consumes its own
output.

\emph{Circularity.} The three circularity patterns of
Remark~\ref{prop:circularity-stress} are blocked by the
non-circularity schema of Proposition~\ref{prop:non-circularity}, which
defines \FATCD{} without reference to a six-slot count, and by the
constructive proof of Theorem~\ref{thm:decomposition-existence}, which
produces a decomposition without assuming scoped exact-six. The absence
of a seventh role is decided by classification under
Definition~\ref{def:residual-scheme} rather than by fiat.

\emph{Channel versus activation.} The conflations of
Remark~\ref{prop:channel-activation-stress} are excluded by
Proposition~\ref{prop:channel-vs-activation} together with the
five-status taxonomy of Definition~\ref{def:channel-status}: a channel
carrying status \(\belowthr\), \(\collapsedbc\), \(\absentrec\), or
\(\outsidescope\) is an inactive channel and is never read as an active
projection, so no description is required to activate all six channels.

\emph{Role alignment.} The three forbidden assignments of
Remark~\ref{prop:role-alignment-stress} are ruled out by the alignment
rules of Proposition~\ref{prop:alignment-rules}, which refuse generic
transition data as \Pthree{}, generic path or derivation data as
\Pfive{}, and generic semantic or interpretation data as \Pone{}, in
accordance with the corrected alignments of
Remark~\ref{rem:three-repaired-alignments}.

\emph{Upper-bound residual.} The regressions of
Remark~\ref{prop:residual-stress} are closed within the covered scope by
Propositions~\ref{prop:probability-closure}--\ref{prop:agency-closure},
which settle the probability, causality, and agency threat closures, and
by Theorem~\ref{thm:upper-bound-classification}, which --- under the
recognition hypothesis (Definition~\ref{def:recognition-hypothesis}) ---
classifies every covered residual feature and delivers the emptiness of
categories (x) and (xi); hence, under that hypothesis, no covered residual
reopens as an unresolved threat or forces a genuine seventh role, while the
three principal threat classes are closed unconditionally.

\emph{Overclaim.} The inadmissible restatements of
Remark~\ref{prop:overclaim-stress} are blocked because the nonclaims of
Remark~\ref{rem:atlas-global-nonclaims} are carried forward by
Proposition~\ref{prop:nonclaim-closure}, and non-uniqueness of the
decomposition is recorded by Proposition~\ref{prop:non-uniqueness}; thus
nonclaim closure survives the exact-six assembly, and no full-algebra,
global-independence, unique-decomposition, empirical-realization, or
broadened-domain reading is admitted.

Each of the six guardrails is therefore satisfied. Dependencies cohere,
the channel-versus-activation distinction holds, nonclaim closure holds,
and the chain admits none of the listed regressions, so the assembled
target is eligible for the terminal statement of
Section~\ref{sec:scoped-exact-six}.
\end{proof}
 \section{The scoped exact-six theorem}\label{sec:scoped-exact-six}




The theorem chain of
Sections~\ref{sec:fatcd-domain}--\ref{sec:decomposition-exhaustion}
and the integrated stress certification of
Section~\ref{sec:integrated-stress} together permit the terminal
statement. The scoped exact-six theorem combines the definition of
\FATCD{} with the lower bounds, the upper-bound classification, and
the decomposition/exhaustion theorems into a single quantified claim
over the domain.

\begin{theorem}[Scoped exact-six]\label{thm:scoped-exact-six}
Let \(D \in \FATCD{}\) satisfy the recognition hypothesis
(Definition~\ref{def:recognition-hypothesis}). Then \(D\) admits a well-formed
decomposition/exhaustion record \(\Dec{D}\) with exactly six typed
closure-role channels:
\begin{itemize}[topsep=2pt,itemsep=1pt]
  \item \Pone{}: update-vs-quotient compatibility / descent obstruction;
  \item \Ptwo{}: macro-specification representability;
  \item \Pthree{}: enriched route/coherence mismatch;
  \item \Pfour{}: stage/order/transition/refinement structure;
  \item \Pfive{}: quotient packaging / quotienting;
  \item \Psix{}: audit/bookkeeping/provenance/replay/checkability.
\end{itemize}
Each channel carries exactly one recorded status from
\(\{\activeproj,\; \belowthr,\; \collapsedbc,\; \absentrec,\;
\outsidescope\}\). Every active projection is supported by threshold
data, a witness schema, a level assignment, and an honest bookkeeping
record. Every residual raw feature of \(D\) is classified, under the
upper-bound classification of
Section~\ref{sec:upper-bound-classification}, as one of:
\begin{itemize}[topsep=2pt,itemsep=1pt]
  \item composite of existing role channels;
  \item presentation variant;
  \item level shift;
  \item activation-threshold refinement;
  \item instrument-visibility artifact;
  \item empirical-bridge annotation;
  \item outside-domain structure;
  \item already-covered bookkeeping or annotation data.
\end{itemize}
Under that hypothesis, no residual feature remains as an unresolved candidate
seventh role or genuine seventh irreducible typed closure role; for the three
principal threat classes --- probability, causality, and agency --- this holds
unconditionally
(Propositions~\ref{prop:probability-closure}--\ref{prop:agency-closure}). We
denote this conclusion by
\[
\ScopedExactSix{\FATCD{}}.
\]\footnote{Tracked in the formalization harness as a conditional schema (\texttt{SixBirds.scoped\_exact\_six}); see Appendix~\ref{app:mechanization}.}
\end{theorem}

\begin{proof}
Fix \(D \in \FATCD{}\) satisfying the recognition hypothesis. By
Theorem~\ref{thm:decomposition-existence},
\(D\) admits a well-formed decomposition/exhaustion record \(\Dec{D}\)
in the sense of Definition~\ref{def:decomposition-record}. By
Theorem~\ref{thm:six-channel-exactness}, the channel sequence of
\(\Dec{D}\) has length six, each channel carries one of the five
statuses of Definition~\ref{def:channel-status}, and every raw
feature of \(D\) is accounted for by an active projection, by a
status record, by a residual classification, or by an auxiliary
loss/lift or visibility annotation. By
Theorem~\ref{thm:upper-bound-classification}, which applies because \(D\)
satisfies the recognition hypothesis, residual features classify under
categories (ii)--(ix) of Definition~\ref{def:residual-scheme}, and categories
(x) and (xi) are empty. Theorem~\ref{thm:six-role-lower-bounds} supplies, for each
channel, a \(D' \in \FATCD{}\) at which that channel is active. The
integrated stress certification of
Proposition~\ref{thm:integrated-stress} discharges the dependency,
circularity, channel-versus-activation, role-alignment, upper-bound
residual, and overclaim families. Hence \(\ScopedExactSix{\FATCD{}}\).
\end{proof}

\subsection{What scoped exact-six means and does not mean}\label{subsec:exact-six-meaning}

Theorem~\ref{thm:scoped-exact-six} has a reading that should be
preserved throughout what follows.

\begin{itemize}[topsep=2pt,itemsep=1pt]
  \item \emph{Six channels, not six active witnesses.} The theorem
    asserts six role channels in \(\Dec{D}\) for every
    \(D \in \FATCD{}\). It does not assert that every \(D\) activates
    all six roles, that every \(D\) carries six active projections
    \(\pi_1(D), \ldots, \pi_6(D)\), or that a single \(D\) realizes
    all six witness schemas of Section~\ref{sec:lower-bounds}.
  \item \emph{Existential, not unique, decomposition.} The theorem
    asserts existence of \(\Dec{D}\); it does not assert uniqueness or
    canonicity. Distinct well-formed records of the same \(D\) may
    exist (Proposition~\ref{prop:non-uniqueness}).
  \item \emph{Conditional residual exhaustion.} The residual
    exhaustion clause is conditional on the recognition hypothesis
    (Definition~\ref{def:recognition-hypothesis}) and restricted to the
    covered scope: the three principal threat classes are closed
    unconditionally, while whether every finite audited feature is recognized
    is the open recognition-completeness question
    (Remark~\ref{rem:recognition-open}). Features requiring non-finite or
    unbookkept semantics are classified as outside-domain structures under
    category~(viii) of Definition~\ref{def:residual-scheme}, not absorbed into
    primitive roles.
  \item \emph{No broader claim.} Theorem~\ref{thm:scoped-exact-six}
    is a statement about \FATCD{}. It does not assert a full
    six-primitives algebra, global primitive independence, empirical
    realization, instrument-family completeness, or coverage of all
    emergence phenomena. See Section~\ref{sec:limitations-nonclaims}.
\end{itemize}

 \section{Discussion}\label{sec:discussion}



The terminal theorem is stated at scoped strength over \FATCD{}.
This section discusses what the statement earns and what it leaves
unstated, how the three parts of
Sections~\ref{sec:primitive-role-architecture},
\ref{sec:admissibility-meta-theory}, and
\ref{sec:fatcd-domain}--\ref{sec:integrated-stress} combine into one
theorem stack, and how the result sits next to earlier Six Birds
literature.

\subsection{The shape of the scoped claim}\label{subsec:discussion-shape}

Each of the three scope features --- domain restriction, existence
rather than uniqueness, residual exhaustion within a covered scope
--- is structural rather than cosmetic. An unrestricted exact-six
theorem would require a domain-free argument for which the present
paper does not claim to have the instruments. A uniqueness or
canonicity claim would require a theory of translation records
identifying distinct well-formed decompositions, which lies outside
the present admissibility regime. Scope-free residual exhaustion
would require scope-free closure of the three high-priority threats
addressed in Section~\ref{sec:upper-bound-classification}. Stating
the theorem at the scope the proof supports keeps the
endogenous-instrument meta-limit (Section~\ref{subsec:meta-limit})
separate from the current claim instead of folded into it. The posture
aligns with the broader view that scientific laws form a patchwork of
domain-local regularities rather than a single global
theory~\citep{Cartwright1999}.

\subsection{The three parts as one theorem stack}\label{subsec:discussion-pillars}

The three parts introduced in Section~\ref{subsec:three-pillars} play
complementary roles. The primitive role architecture
(Section~\ref{sec:primitive-role-architecture}) fixes the vocabulary:
six role meanings, the behavioral, verification, and
structural-categorical level pattern, and the local boundary
distinctions. This vocabulary alone does not certify any claim,
since it does not specify which role talk is admissible. The
admissibility and meta-theory regime
(Section~\ref{sec:admissibility-meta-theory}) supplies that
certification through honest bookkeeping, typed non-collapse, the
forgetting and selected-lift architecture, activation thresholds,
instrument-relative visibility, and empirical-bridge admissibility.
Admissibility alone, however, yields no terminal claim; it only
regulates what may be asserted. The exact-six theorem chain
(Sections~\ref{sec:fatcd-domain}--\ref{sec:integrated-stress})
supplies the quantification: the \FATCD{} domain, the six
lower-bound witnesses, the upper-bound classification, the
decomposition/exhaustion theorem, and the integrated stress
certification together deliver Theorem~\ref{thm:scoped-exact-six}.
Vocabulary without admissibility would be undisciplined; admissibility
without a theorem chain would be inconclusive; a theorem chain without
the first two would lack a subject.

\subsection{Channel versus activation}\label{subsec:discussion-channel-vs-activation}

The channel-versus-activation discipline of
Proposition~\ref{prop:channel-vs-activation}, enforced throughout
Sections~\ref{sec:fatcd-domain}--\ref{sec:scoped-exact-six},
separates two statements that are easily conflated. The structural
statement is that every admissible description carries six typed
closure-role slots. The activation statement is that every admissible
description realizes all six slots with active projections. The
activation statement is too strong: it fails on descriptions that
carry no route grammar and therefore have no active \Pthree{}
projection, and on descriptions whose raw material carries no
\Psix{}-aligned event/provenance/replay/check witness cluster and
therefore have no active \Psix{} projection. The structural
statement, which is what we prove, is compatible with any activation
profile. Inactive channels carry one of the statuses
\(\belowthr{}\), \(\collapsedbc{}\), \(\absentrec{}\), or
\(\outsidescope{}\), describing the way a role fails to activate.
Separating the two statements is what lets the terminal theorem hold
of every \(D \in \FATCD{}\) without requiring activation of all six
roles in every description; it also lets the lower-bound theorem of
Section~\ref{sec:lower-bounds} be stated existentially over the
domain rather than uniformly. An informal reading that elides the
distinction silently recovers the stronger activation claim, which we
do not prove.

\subsection{Instrument-relative methodology}\label{subsec:discussion-instrument-relative}

A methodological choice underlying
Section~\ref{sec:admissibility-meta-theory} deserves explicit
statement as a contribution. Instruments are treated as mathematical
probes rather than as neutral tools, in the spirit of ontic structural
realism~\citep{LadymanRoss2007} and the sheaf-theoretic treatment of
contextuality~\citep{AbramskyBrandenburger2011}. The technical form
is Definition~\ref{prop:visibility}; the associated nonclaim is that
no result here rests on an instrument-free truth.

The motivation is reflexive. The subject matter --- finite audited
typed closure descriptions and the six typed closure roles that
project them --- is the same descriptional machinery from which any
instrument for detecting closure structure must itself be built.
Operator-presentation, quotient/descent, coherence/route,
constructive/proof-theoretic, and information/regime instruments are
not vantage points outside the domain: each is a closure structure
probing another closure structure. A claim asserted under one such
instrument while silently suppressing what another would have
revealed would privilege that instrument without bookkeeping for the
privilege --- the overread that Definition~\ref{prop:visibility}
blocks.

The consequence is epistemic. Every admissible claim carries an
explicit visibility record: its visible content, its suppressed
content, the instruments under which visibility and suppression are
asserted, and the resulting instrument-relative claim strength
(Definition~\ref{prop:visibility}). Triangulating a role projection
across several instruments is a disciplined way to strengthen an
assertion; reading a projection off a single instrument and declaring
its content canonical is not admissible. The discipline is one of
epistemic hygiene, in the spirit of Bohr's complementarity
doctrine~\citep{Bohr1958}: we do not claim that instruments are
incompatible in principle, only that they are not interchangeable
without explicit visibility bookkeeping and, where contexts shift,
explicit translation records
(Definitions~\ref{prop:visibility} and~\ref{prop:level-profile}).

This posture also explains why the endogenous-instrument meta-limit
(Sections~\ref{subsec:meta-limit} and~\ref{subsec:future-meta-limit})
is deferred to future work. Taking instrument-relativity seriously
from the outset lets the meta-limit be stated as its own question ---
whether the accepted internal instrument family can certify its own
soundness and completeness --- a question made statable, but not
settled, by the stance adopted here.

\subsection{Relation to earlier Six Birds results}\label{subsec:discussion-earlier-results}

Six Birds Foundations~I~\citep{Tsiokos2026SixBirdsFoundations}
introduced the theory-package formalism
\(\mathcal{T} = (Z, f, \Sigma_f, E, \mathcal{A})\) and derived the
\Pone{}--\Psix{} vocabulary under minimal hypotheses of composability
with limited interface access. The present paper does not re-derive
that vocabulary; it takes the six roles as given and asks a
different question: whether, within a specific admissibility regime
and over a specific class of admissible descriptions, the six roles
are exactly enough. The Foundations theorem gives the canonical
unavoidability side of the picture; the terminal theorem here gives
the scoped exhaustion side. The two statements address different
questions and address them at different levels; combining them does
not yield a stronger claim, and neither theorem follows from the
other.

The theorem chain also draws on a nearby
result. \citet{Tsiokos2026SixBirdsProtocolTrap}, complementing the
no-go analysis of \citet{Tsiokos2026SixBirdsNoGo}, established that
protocol holonomy is not a directionality certificate:
noncommutativity of update sequences can be nonzero while
path-reversal asymmetry remains zero, so route mismatch is a
diagnostic, not an arrow of time. That observation is load-bearing
for the role-projection alignment rules of
Section~\ref{sec:upper-bound-classification}, where \Pthree{}
route/coherence mismatch is kept structurally distinct from the
affinity-carrying face of \Psix{}. The next section records the
scope boundaries and nonclaims of the present theorem stack.
 \section{Limitations and nonclaims}\label{sec:limitations-nonclaims}



This section collects the scope restrictions of
Theorem~\ref{thm:scoped-exact-six} and the global nonclaims
preserved throughout.

\subsection{Scope of the terminal theorem}\label{subsec:terminal-scope}

Theorem~\ref{thm:scoped-exact-six} quantifies over \(D \in \FATCD{}\),
where \FATCD{} is the finite audited typed closure description domain
of Definition~\ref{def:fatcd}. The theorem asserts:
\begin{itemize}[topsep=2pt,itemsep=1pt]
  \item existence of a decomposition/exhaustion record \(\Dec{D}\)
    with six role channels and one status per channel;
  \item residual-feature exhaustion under the upper-bound classification,
    conditional on the recognition hypothesis
    (Definition~\ref{def:recognition-hypothesis}) and within the covered scope,
    so that no seventh irreducible role is required there (the three principal
    threat classes are closed unconditionally);
  \item activation of any \(\pi_i(D)\) only under the threshold,
    witness, level, and audit data of
    Definition~\ref{def:threshold-attachment}.
\end{itemize}

It does not assert:
\begin{itemize}[topsep=2pt,itemsep=1pt]
  \item decomposition uniqueness or canonicity;
  \item six active projections, or six active nontrivial witnesses, in
    every description;
  \item scope-free residual exhaustion, or any claim about features
    outside the covered scope;
  \item that the six roles are recognition-complete for \FATCD{}, i.e.\ that
    every finite audited feature reduces to them;
  \item unrestricted exact-six, or coverage of all emergence
    phenomena;
  \item that broader domains require no further work.
\end{itemize}

\subsection{Global nonclaims}\label{subsec:global-nonclaims}

Each nonclaim below is preserved throughout the paper and is enforced
locally by the theorem-chain result that would otherwise overread it.

\begin{itemize}[topsep=2pt,itemsep=1pt]
  \item \emph{No full six-primitives algebra.} Scoped exact-six is
    proved over \FATCD{}; it is not claimed that \Pone{}--\Psix{}
    generate a global algebra of closure operations
    (Remark~\ref{rem:atlas-global-nonclaims},
    Remark~\ref{prop:overclaim-stress}).
  \item \emph{No global primitive-independence theorem.} Typed
    non-collapse (Proposition~\ref{prop:typed-non-collapse}) is
    strictly weaker than global independence and is the sole claim at
    that level.
  \item \emph{No empirical realization.} Empirical-bridge admissibility
    (Definition~\ref{prop:empirical-bridge}) admits
    threshold-dependent, level-indexed, instrument-visible substrate
    evidence; it does not assert that any \(D \in \FATCD{}\) is itself
    a substrate observation.
  \item \emph{No unique selected lift.} Every lift is recorded as
    selected and non-unique
    (Definition~\ref{prop:forgetting-lifts}).
  \item \emph{No lower-to-higher determinacy.} Behavioral or
    verification content does not determine structural-categorical
    content.
  \item \emph{No instrument-free formulation.} Admissible claims carry
    an instrument-relative visibility record
    (Definition~\ref{prop:visibility}).
  \item \emph{No internal instrument-family completeness.} The
    accepted instrument family is not asserted to be sufficient for
    every exact-six judgment formulable in the same regime.
  \item \emph{No unrestricted all-emergence coverage.}
    Theorem~\ref{thm:scoped-exact-six} is scoped to \FATCD{}.
  \item \emph{No surrogate terminal evidence.} The proof chain does
    not invoke any surrogate for the terminal scoped exact-six
    theorem.
  \item \emph{No uniqueness of primitive presentations.} Distinct
    admissible presentations of the same role may coexist.
  \item \emph{Cross-level comparison requires translation.} Direct
    comparison of claims at different levels is inadmissible without a
    translation record (Definition~\ref{prop:level-profile}).
\end{itemize}

\subsection{The endogenous-instrument meta-limit}\label{subsec:meta-limit}

An adjacent meta-theoretic question, deferred to future work, is
whether the accepted internal instrument family can itself certify
both its soundness and its completeness for exact-six judgments
formulable inside the same regime. A plausible sharp formulation is
that no internal argument should be expected to accomplish all three
of the following at once:
\begin{enumerate}[topsep=2pt,itemsep=1pt]
  \item establish the scoped exact-six claim;
  \item certify the soundness of the accepted instrument family;
  \item certify the completeness of that family for exact-six
    judgments expressible inside the same regime.
\end{enumerate}
Such self-certification limits are of a piece with the classical
G{\"o}del--L{\"o}b tradition on internal consistency and provability
statements~\citep{Godel1931,Lob1955}. This question is not part of
the proof of Theorem~\ref{thm:scoped-exact-six} and does not weaken
or retract the theorem's scoped conclusion. It is taken up as a
follow-on direction in Section~\ref{sec:future-work}.
 \section{Future work}\label{sec:future-work}



Theorem~\ref{thm:scoped-exact-six} and the limitations of
Section~\ref{sec:limitations-nonclaims} suggest several follow-on
research directions. Each subsection below states a direction, not a
commitment.

\subsection{Endogenous-instrument meta-limit}\label{subsec:future-meta-limit}

A formalization of the meta-limit described in
Section~\ref{subsec:meta-limit} would target a three-way
impossibility statement: no internal argument can simultaneously
establish the scoped exact-six claim, certify the soundness of the
accepted instrument family, and certify its completeness for
exact-six judgments formulable inside the same regime.
Theorem~\ref{thm:scoped-exact-six} would remain intact under such a
result; the additional content would be a precise account of
self-certification boundaries. A closely related open question is
\emph{recognition-completeness}: whether every finite audited feature of a
description in \FATCD{} reduces to the six role meanings --- the recognition
hypothesis on which Theorem~\ref{thm:upper-bound-classification} is conditional
(Remark~\ref{rem:recognition-open}). Candidate witnesses against it --- finite
conservation laws and emergent composite obstructions --- would, if irreducible,
populate the genuine-seventh-role category; establishing or refuting their
reducibility is left to future work.

\subsection{Broader-domain exact-six}\label{subsec:future-broader-domain}

Candidate extensions beyond \FATCD{} include admissibility regimes
that weaken finiteness, closure, or audit
(Section~\ref{sec:fatcd-domain}), and regimes that admit annotation
kinds beyond the six \(\mathrm{Ann}_*\) categories of
Definition~\ref{def:annotations}. Any such extension would require
fresh upper-bound and decomposition/exhaustion analysis, since the
alignment rules and threat closures are not stated at scope-free
strength.

\subsection{Uniqueness or canonical decomposition}\label{subsec:future-uniqueness}

Theorem~\ref{thm:scoped-exact-six} is existential. A uniqueness study
would investigate when distinct well-formed decomposition/exhaustion
records of the same \(D \in \FATCD{}\) are equivalent under explicit
translation records, or when a canonical representative can be
selected. The channel-versus-activation discipline
(Proposition~\ref{prop:channel-vs-activation}) is compatible with
such a study; no uniqueness claim is part of the present theorem.

\subsection{Full six-primitives algebra}\label{subsec:future-algebra}

The generator-and-relation structure of \Pone{}--\Psix{} beyond
typed non-collapse remains open. A full algebra would require
scope-free independence data, not just the local boundary
distinctions of Table~\ref{tab:boundary-matrix}, and would likely
require a categorical apparatus richer than the raw grammar of
Section~\ref{subsec:raw-grammar}.

\subsection{Empirical substrate mapping}\label{subsec:future-substrate}

The empirical-bridge admissibility regime
(Definition~\ref{prop:empirical-bridge}) can be connected to
concrete substrate-level observables. The present paper treats
empirical bridges as threshold-dependent, level-indexed, and
instrument-visible annotations; later work can couple this to
substrate-specific measurement theory while preserving the nonclaim
of empirical realization.

\subsection{Mechanization}\label{subsec:future-mechanization}

A Lean~4 development accompanies the present paper;
Appendix~\ref{app:mechanization} records, claim by claim, how far each definition
and theorem is currently mechanized. Strengthening that mechanization toward full
proofs --- in particular a Lean account of the boundary distinctions, the
decomposition-existence construction, and the recognition hypothesis underlying
the upper-bound classification --- is a natural direction, building on the
closure-related infrastructure already mechanized in the Foundations
program~\citep{Tsiokos2026SixBirdsFoundations,DeMouraUllrich2021}.

\subsection{Cross-instrument comparison}\label{subsec:future-cross-instrument}

Decomposition/exhaustion records of the same \(D \in \FATCD{}\) can
be compared under different admissible instrument assignments.
Definition~\ref{prop:visibility} admits visibility records that
suppress different portions of a claim under different instruments; a
cross-instrument study could make the instrument-dependence of
\(\Dec{D}\) explicit and sharpen the notion of instrument-relative
claim strength.
 
\appendix
\section{Provenance and evidence hygiene}\label{app:provenance}



The development of the theorem chain of
Sections~\ref{sec:fatcd-domain}--\ref{sec:integrated-stress} was assisted by
automated drafting and proof-search tooling, the Aristotle
system~\citep{Achim2025Aristotle}. We record this for transparency and to fix the
division of evidential labor precisely. Such tooling produces \emph{candidate}
definitions, statements, and arguments; it carries no evidential weight of its
own, and nothing in this paper rests on it.

The mathematical content is established solely by the proofs given in the body.
The machine-\emph{checked} component of the work --- which declarations are
realized in Lean, and at what fidelity --- is documented separately in
Appendix~\ref{app:mechanization}; the automated tooling named above is distinct
from that mechanization and assisted drafting only. No raw output of that tooling
is offered as mathematical evidence anywhere in the manuscript: where a section
draws on machine assistance, the load-bearing content is the section's own proof,
and the recorded mechanization is exactly as described in
Appendix~\ref{app:mechanization}.

\subsection{Code availability}\label{app:code-availability}

The \LaTeX{} source of this manuscript and the accompanying Lean mechanization
artifacts are available at
\url{https://github.com/ioannist/six-birds-foundations-ii}.
 \section{Mechanization: declarations, coverage, and trust base}\label{app:mechanization}

The theorem chain of this paper is accompanied by a Lean~4 development, the
integrated \texttt{SixBirds} project, available at the code repository of
Appendix~\ref{app:code-availability}. This appendix is the canonical disclosure
venue for that development. For every numbered definition and claim it records
the Lean declaration that realizes it and the \emph{fidelity} of that
realization, so the per-theorem footnotes in the body can be read against a
single source of record. The mapping is maintained against a statements-of-record
file in the repository and reproduced here; coverage is audited per declaration
against that file.

\subsection{What the footnotes mean}\label{app:fidelity-legend}

Each body theorem and proposition carries exactly one footnote naming its Lean
declaration and its fidelity. We distinguish five fidelities:
\begin{itemize}[topsep=2pt,itemsep=1pt]
  \item \emph{verified} --- a faithful Lean derivation of the statement from
    non-trivial hypotheses;
  \item \emph{narrowed / parametric} --- a genuine Lean proof, but over a single
    concrete witness, a parametric record, or a narrowed input space rather than
    the full generality of the paper statement;
  \item \emph{schema (projection)} --- the Lean declaration projects a field of,
    or restates, a proof-carrying record; it repackages rather than derives, and
    is therefore \emph{not} described as a verified proof;
  \item \emph{trivial} --- the Lean statement reduces to \texttt{True} by
    construction, or ranges over a hand-built total classifier; it records a
    bookkeeping fact, not a derivation;
  \item \emph{typed mirror} --- for definitions, a Lean structure or inductive
    mirroring the paper definition.
\end{itemize}
We state plainly that \textbf{no claim in the present development is a faithful
derivation in the first sense}. The Lean track is a typed harness: it fixes the
objects as typed records, exposes the channel, status, and residual machinery,
and checks that malformed and misaligned cases are rejected (the semantic
regression tests of the \texttt{SixBirds.SemanticTests} module); but the
load-bearing classification and closure statements are projections, schemas, or
trivially-true bookkeeping over those records. The footnotes are worded
accordingly. A reader who wants the mathematics should read the body proofs; the
Lean track certifies typed well-formedness and case rejection, not the theorems.

\subsection{The recognition hypothesis and the upper-bound classifier}\label{app:recognition}

The honest status of the central no-seventh-role result
(Theorem~\ref{thm:upper-bound-classification}) is legible in the Lean. The
classifier \texttt{SixBirds.classifyResidual} is a total function on the finite
raw-record kinds, every branch of which lands in a covered class; the
declarations \texttt{upper\_bound\_classification} and the three closure lemmas
are \texttt{True} by construction over it. This is exactly why the body states
the result conditionally, under the recognition hypothesis
(Definition~\ref{def:recognition-hypothesis}): the Lean classifier \emph{is} that
hypothesis, encoded as a conditional schema rather than discharged. The
paper-side recognition hypothesis has no independent Lean correspondent, and the
open recognition-completeness question (Remark~\ref{rem:recognition-open}) is
precisely the gap between the classifier's by-construction coverage and a derived
exhaustiveness.

\subsection{Declaration and coverage table}\label{app:table}

The table maps each paper label to its Lean declaration in the \texttt{SixBirds}
namespace and its audited coverage. All declarations are \texttt{sorry}-free and
axiom-clean relative to the trust base of \S\ref{app:trust-base}; the coverage
column records fidelity in the sense of \S\ref{app:fidelity-legend}, not mere
presence.

{\footnotesize
\begin{longtable}{@{}c >{\raggedright\arraybackslash}p{0.27\linewidth} >{\raggedright\arraybackslash}p{0.34\linewidth} >{\raggedright\arraybackslash}p{0.20\linewidth}@{}}
\toprule
\S & paper label & Lean declaration & coverage \\
\midrule
\endfirsthead
\toprule
\S & paper label & Lean declaration & coverage \\
\midrule
\endhead
\bottomrule
\endfoot
§2 & \texttt{def:\allowbreak primitive-\allowbreak roles} & \texttt{SixBirds.\allowbreak Role} & typed mirror (def.) \\
§2 & \texttt{def:\allowbreak level-\allowbreak trichotomy} & \texttt{SixBirds.\allowbreak Level} & typed mirror (def.) \\
§2 & \texttt{def:\allowbreak selected-\allowbreak lift-\allowbreak atlas} & \texttt{SixBirds.\allowbreak SelectedLiftAtlas} & typed mirror (def.) \\
§2 & \texttt{prop:\allowbreak boundary-\allowbreak distinctions} & \texttt{SixBirds.\allowbreak boundary\_\allowbreak distinctions} & schema (projection) \\
§3 & \texttt{prop:\allowbreak honest-\allowbreak bookkeeping} & \texttt{SixBirds.\allowbreak honest\_\allowbreak bookkeeping\_\allowbreak admissibility} & schema (projection) \\
§3 & \texttt{prop:\allowbreak typed-\allowbreak non-\allowbreak collapse} & \texttt{SixBirds.\allowbreak typed\_\allowbreak non\_\allowbreak collapse} & schema (projection) \\
§3 & \texttt{prop:\allowbreak level-\allowbreak profile} & \texttt{SixBirds.\allowbreak level\_\allowbreak profile} & trivial \\
§3 & \texttt{prop:\allowbreak forgetting-\allowbreak lifts} & \texttt{SixBirds.\allowbreak forgetting\_\allowbreak lifts} & schema (projection) \\
§3 & \texttt{prop:\allowbreak activation-\allowbreak thresholds} & \texttt{SixBirds.\allowbreak activation\_\allowbreak thresholds} & schema (projection) \\
§3 & \texttt{prop:\allowbreak visibility} & \texttt{SixBirds.\allowbreak visibility} & schema (projection) \\
§3 & \texttt{prop:\allowbreak empirical-\allowbreak bridge} & \texttt{SixBirds.\allowbreak empirical\_\allowbreak bridge} & schema (projection) \\
§3 & \texttt{prop:\allowbreak nonclaim-\allowbreak closure} & \texttt{SixBirds.\allowbreak nonclaim\_\allowbreak closure} & schema (projection) \\
§4 & \texttt{def:\allowbreak raw-\allowbreak record} & \texttt{SixBirds.\allowbreak RawRecord} & typed mirror (def.) \\
§4 & \texttt{def:\allowbreak annotations} & \texttt{SixBirds.\allowbreak AnnotationKind} & typed mirror (def.) \\
§4 & \texttt{def:\allowbreak finite} & \texttt{SixBirds.\allowbreak FiniteDescription} & typed mirror (def.) \\
§4 & \texttt{def:\allowbreak closed} & \texttt{SixBirds.\allowbreak ClosedDescription} & typed mirror (def.) \\
§4 & \texttt{def:\allowbreak audited} & \texttt{SixBirds.\allowbreak AuditedDescription} & typed mirror (def.) \\
§4 & \texttt{def:\allowbreak fatcd} & \texttt{SixBirds.\allowbreak FATCD} & typed mirror (def.) \\
§4 & \texttt{def:\allowbreak threshold-\allowbreak attachment} & \texttt{SixBirds.\allowbreak ThresholdAttachment} & typed mirror (def.) \\
§4 & \texttt{def:\allowbreak role-\allowbreak projection} & \texttt{SixBirds.\allowbreak DerivedRoleProjection} & typed mirror (def.) \\
§4 & \texttt{def:\allowbreak residual-\allowbreak scheme} & \texttt{SixBirds.\allowbreak ResidualClassificationScheme} & typed mirror (def.) \\
§4 & \texttt{def:\allowbreak scoped-\allowbreak exact-\allowbreak six-\allowbreak question} & \texttt{SixBirds.\allowbreak ScopedExactSixQuestion} & typed mirror (def.) \\
§4 & \texttt{prop:\allowbreak non-\allowbreak circularity} & \texttt{SixBirds.\allowbreak non\_\allowbreak circularity} & schema (projection) \\
§4 & \texttt{prop:\allowbreak non-\allowbreak overbreadth} & \texttt{SixBirds.\allowbreak non\_\allowbreak overbreadth} & schema (projection) \\
§5 & \texttt{def:\allowbreak common-\allowbreak witness-\allowbreak data} & \texttt{SixBirds.\allowbreak CommonWitnessData} & typed mirror (def.) \\
§5 & \texttt{thm:\allowbreak six-\allowbreak role-\allowbreak lower-\allowbreak bounds} & \texttt{SixBirds.\allowbreak six\_\allowbreak role\_\allowbreak lower\_\allowbreak bounds} & schema (projection) \\
§5 & \texttt{prop:\allowbreak p1-\allowbreak witness} & \texttt{SixBirds.\allowbreak p1\_\allowbreak witness} & narrowed / parametric \\
§5 & \texttt{prop:\allowbreak p2-\allowbreak witness} & \texttt{SixBirds.\allowbreak p2\_\allowbreak witness} & narrowed / parametric \\
§5 & \texttt{prop:\allowbreak p3-\allowbreak witness} & \texttt{SixBirds.\allowbreak p3\_\allowbreak witness} & narrowed / parametric \\
§5 & \texttt{prop:\allowbreak p4-\allowbreak witness} & \texttt{SixBirds.\allowbreak p4\_\allowbreak witness} & narrowed / parametric \\
§5 & \texttt{prop:\allowbreak p5-\allowbreak witness} & \texttt{SixBirds.\allowbreak p5\_\allowbreak witness} & narrowed / parametric \\
§5 & \texttt{prop:\allowbreak p6-\allowbreak witness} & \texttt{SixBirds.\allowbreak p6\_\allowbreak witness} & narrowed / parametric \\
§6 & \texttt{prop:\allowbreak alignment-\allowbreak rules} & \texttt{SixBirds.\allowbreak alignment\_\allowbreak rules} & schema (projection) \\
§6 & \texttt{prop:\allowbreak probability-\allowbreak closure} & \texttt{SixBirds.\allowbreak probability\_\allowbreak closure} & trivial \\
§6 & \texttt{prop:\allowbreak causality-\allowbreak closure} & \texttt{SixBirds.\allowbreak causality\_\allowbreak closure} & trivial \\
§6 & \texttt{prop:\allowbreak agency-\allowbreak closure} & \texttt{SixBirds.\allowbreak agency\_\allowbreak closure} & trivial \\
§6 & \texttt{thm:\allowbreak upper-\allowbreak bound-\allowbreak classification} & \texttt{SixBirds.\allowbreak upper\_\allowbreak bound\_\allowbreak classification} & trivial \\
§7 & \texttt{def:\allowbreak channel-\allowbreak status} & \texttt{SixBirds.\allowbreak ChannelStatus} & typed mirror (def.) \\
§7 & \texttt{def:\allowbreak decomposition-\allowbreak record} & \texttt{SixBirds.\allowbreak DecompositionRecord} & typed mirror (def.) \\
§7 & \texttt{thm:\allowbreak decomposition-\allowbreak existence} & \texttt{SixBirds.\allowbreak decomposition\_\allowbreak existence} & schema (projection) \\
§7 & \texttt{thm:\allowbreak six-\allowbreak channel-\allowbreak exactness} & \texttt{SixBirds.\allowbreak six\_\allowbreak channel\_\allowbreak exactness} & schema (projection) \\
§7 & \texttt{prop:\allowbreak channel-\allowbreak vs-\allowbreak activation} & \texttt{SixBirds.\allowbreak channel\_\allowbreak vs\_\allowbreak activation} & schema (projection) \\
§7 & \texttt{prop:\allowbreak non-\allowbreak uniqueness} & \texttt{SixBirds.\allowbreak non\_\allowbreak uniqueness} & narrowed / parametric \\
§8 & \texttt{prop:\allowbreak dependency-\allowbreak stress} & \texttt{SixBirds.\allowbreak dependency\_\allowbreak stress} & schema (projection) \\
§8 & \texttt{prop:\allowbreak circularity-\allowbreak stress} & \texttt{SixBirds.\allowbreak circularity\_\allowbreak stress} & schema (projection) \\
§8 & \texttt{prop:\allowbreak channel-\allowbreak activation-\allowbreak stress} & \texttt{SixBirds.\allowbreak channel\_\allowbreak activation\_\allowbreak stress} & narrowed / parametric \\
§8 & \texttt{prop:\allowbreak role-\allowbreak alignment-\allowbreak stress} & \texttt{SixBirds.\allowbreak role\_\allowbreak alignment\_\allowbreak stress} & narrowed / parametric \\
§8 & \texttt{prop:\allowbreak residual-\allowbreak stress} & \texttt{SixBirds.\allowbreak residual\_\allowbreak stress} & trivial \\
§8 & \texttt{prop:\allowbreak overclaim-\allowbreak stress} & \texttt{SixBirds.\allowbreak overclaim\_\allowbreak stress} & schema (projection) \\
§8 & \texttt{thm:\allowbreak integrated-\allowbreak stress} & \texttt{SixBirds.\allowbreak integrated\_\allowbreak stress} & schema (projection) \\
§9 & \texttt{thm:\allowbreak scoped-\allowbreak exact-\allowbreak six} & \texttt{SixBirds.\allowbreak scoped\_\allowbreak exact\_\allowbreak six} & schema (projection) \\
\bottomrule
\end{longtable}
}

\subsection{Trust base and axiom audit}\label{app:trust-base}

The dependency audit accepts only the standard Lean and Mathlib trust base:
\texttt{Classical.choice}, \texttt{propext}, \texttt{Quot.sound}, and
\texttt{funext} (together with \texttt{choice}). No project-local axiom is
admitted for any declaration recorded as a theorem. The validation harness
(\texttt{scripts/validate\_lean.sh}) rejects project-local axioms and
\texttt{sorry}, and runs a semantic audit that rejects \texttt{:= True} and
\texttt{:= False} definition stubs and all-purpose residual classifiers. The
trust-base check certifies that the typed harness is axiom-clean relative to that
base; it does not, and is not intended to, upgrade a schema or a trivially-true
statement to a derivation. The fidelity column of \S\ref{app:table} is the
governing record of what each declaration establishes.
 


\begin{landscape}
\begin{table}
\caption{Local boundary distinctions. Each row records the distinction
rule between two roles, a permitted interaction, and the collapse that
is forbidden under scoped bookkeeping.}
\label{tab:boundary-matrix}
\centering
\footnotesize
\renewcommand{\arraystretch}{1.15}
\begin{tabular}{@{}lp{0.34\linewidth}p{0.24\linewidth}p{0.24\linewidth}@{}}
\toprule
pair & distinction rule & permitted interaction & forbidden collapse \\
\midrule
\Pone{}/\Pfive{} & \Pfive{} supplies quotient packaging; \Pone{} tests selected update compatibility with quotient packaging & \Pone{} uses \Pfive{} when checking descent or obstruction & packaging alone is not \Pone{} \\
\Pone{}/\Pfour{} & \Pfour{} supplies stage/refinement structure; \Pone{} may vary by stage & stage-indexed descent or obstruction & stage index alone is not \Pone{} \\
\Pone{}/\Ptwo{} & \Pone{} is update compatibility; \Ptwo{} is macro-specification representability & representability may be studied over quotient structures affected by descent & representability failure is not automatically descent obstruction \\
\Pone{}/\Pthree{} & \Pone{} is failure or success of a descended route; \Pthree{} compares already-defined routes & \Pone{} may block a route before \Pthree{} comparison & route undefinedness is not route mismatch \\
\Pone{}/\Psix{} & \Psix{} audits a \Pone{} event; \Pone{} is the event content & audited descent success or obstruction & audit record is not descent obstruction \\
\Ptwo{}/\Pfive{} & \Pfive{} is quotient packaging; \Ptwo{} requires macro-specification representability beyond quotient equality & \Ptwo{} can collapse to \Pfive{} when specifications are just quotient classes; nontrivial \Ptwo{} uses \Pfive{} & nontrivial \Ptwo{} is not mere packaging \\
\Ptwo{}/\Pfour{} & \Pfour{} supplies stages; \Ptwo{} representability may vary by stage & stage-indexed representability & stage variation alone is not \Ptwo{} \\
\Ptwo{}/\Pthree{} & \Ptwo{} concerns representability of specifications; \Pthree{} concerns route disagreement & route grammar may be built over represented macro structure & non-representability is not route mismatch \\
\Ptwo{}/\Psix{} & \Psix{} audits representability claims; \Ptwo{} is the representability content & audited representability or non-representability & audit trail is not representability \\
\Pthree{}/\Pfour{} & \Pfour{} supplies stage/regime structure; \Pthree{} route mismatch may be stage-indexed & route systems vary over valid stages & stage-indexed route mismatch is not \Pfour{} itself \\
\Pthree{}/\Pfive{} & \Pfive{} supplies quotient objects; \Pthree{} compares defined routes over objects & routes over quotient objects & quotient object is not route mismatch \\
\Pthree{}/\Psix{} & \Psix{} audits route comparisons; \Pthree{} is the route mismatch content & audited route mismatch & audit record is not route mismatch \\
\Pfour{}/\Pfive{} & \Pfour{} supplies stage/refinement structure; \Pfive{} supplies packaging at a stage & stage-indexed packaging & packaging is not stage structure \\
\Pfour{}/\Psix{} & \Pfour{} supplies staged structure; \Psix{} audits staged histories or transitions & audited stage transitions & audit record is not staging \\
\Pfive{}/\Psix{} & \Pfive{} supplies packaging; \Psix{} audits packaging events & audited packaging & audit record is not packaging map \\
\bottomrule
\end{tabular}
\end{table}
\end{landscape}
 
% Bibliography embedded from BibTeX .bbl output (no external .bib needed)
\begin{thebibliography}{31}
\providecommand{\natexlab}[1]{#1}
\providecommand{\url}[1]{\texttt{#1}}
\expandafter\ifx\csname urlstyle\endcsname\relax
  \providecommand{\doi}[1]{doi: #1}\else
  \providecommand{\doi}{doi: \begingroup \urlstyle{rm}\Url}\fi

\bibitem[Abramsky and Brandenburger(2011)]{AbramskyBrandenburger2011}
Samson Abramsky and Adam Brandenburger.
\newblock The sheaf-theoretic structure of non-locality and contextuality.
\newblock \emph{New Journal of Physics}, 13\penalty0 (11):\penalty0 113036,
  2011.
\newblock \doi{10.1088/1367-2630/13/11/113036}.

\bibitem[Achim et~al.(2025)Achim, Best, Bietti, Der, F{\'e}d{\'e}rico, Gukov,
  Halpern-Leistner, Henningsgard, Kudryashov, Meiburg, Michelsen, Patterson,
  Rodriguez, Scharff, Shanker, Sicca, Sowrirajan, Swope, Tamas, Tenev, Thomm,
  Williams, and Wu]{Achim2025Aristotle}
Tudor Achim, Alex Best, Alberto Bietti, Kevin Der, Math{\"i}s F{\'e}d{\'e}rico,
  Sergei Gukov, Daniel Halpern-Leistner, Kirsten Henningsgard, Yury Kudryashov,
  Alexander Meiburg, Martin Michelsen, Riley Patterson, Eric Rodriguez, Laura
  Scharff, Vikram Shanker, Vladmir Sicca, Hari Sowrirajan, Aidan Swope, Matyas
  Tamas, Vlad Tenev, Jonathan Thomm, Harold Williams, and Lawrence Wu.
\newblock Aristotle: {IMO}-level automated theorem proving, 2025.

\bibitem[Anderson(1972)]{Anderson1972}
P.~W. Anderson.
\newblock More is different.
\newblock \emph{Science}, 177\penalty0 (4047):\penalty0 393--396, 1972.
\newblock \doi{10.1126/science.177.4047.393}.

\bibitem[Baader and Nipkow(1998)]{BaaderNipkow1998}
Franz Baader and Tobias Nipkow.
\newblock \emph{Term Rewriting and All That}.
\newblock Cambridge University Press, 1998.
\newblock \doi{10.1017/CBO9781139172752}.

\bibitem[Back and von Wright(1998)]{BackVonWright1998}
Ralph-Johan Back and Joakim von Wright.
\newblock \emph{Refinement Calculus: A Systematic Introduction}.
\newblock Texts in Computer Science. Springer, 1998.
\newblock \doi{10.1007/978-1-4612-1674-2}.

\bibitem[Bohr(1958)]{Bohr1958}
Niels Bohr.
\newblock \emph{Atomic Physics and Human Knowledge}.
\newblock John Wiley \& Sons, 1958.

\bibitem[Bratman(1987)]{Bratman1987}
Michael~E. Bratman.
\newblock \emph{Intention, Plans, and Practical Reason}.
\newblock Harvard University Press, 1987.
\newblock ISBN 0674458184.

\bibitem[Cardelli and Mitchell(1991)]{CardelliMitchell1991}
Luca Cardelli and John~C. Mitchell.
\newblock Operations on records.
\newblock \emph{Mathematical Structures in Computer Science}, 1\penalty0
  (1):\penalty0 3--48, 1991.
\newblock \doi{10.1017/S0960129500000049}.

\bibitem[Cartwright(1999)]{Cartwright1999}
Nancy Cartwright.
\newblock \emph{The Dappled World: A Study of the Boundaries of Science}.
\newblock Cambridge University Press, 1999.
\newblock \doi{10.1017/CBO9781139167093}.

\bibitem[Davey and Priestley(2002)]{DaveyPriestley2002}
B.~A. Davey and H.~A. Priestley.
\newblock \emph{Introduction to Lattices and Order}.
\newblock Cambridge University Press, 2 edition, 2002.
\newblock \doi{10.1017/CBO9780511809088}.

\bibitem[{de Moura} and Ullrich(2021)]{DeMouraUllrich2021}
Leonardo {de Moura} and Sebastian Ullrich.
\newblock The {L}ean 4 theorem prover and programming language.
\newblock In Andr{\'e} Platzer and Geoff Sutcliffe, editors, \emph{Automated
  Deduction --- CADE 28}, volume 12699 of \emph{Lecture Notes in Computer
  Science}, pages 625--635. Springer International Publishing, 2021.
\newblock \doi{10.1007/978-3-030-79876-5_37}.

\bibitem[G{\"o}del(1931)]{Godel1931}
Kurt G{\"o}del.
\newblock {\"U}ber formal unentscheidbare {S}{\"a}tze der {P}rincipia
  {M}athematica und verwandter {S}ysteme {I}.
\newblock \emph{Monatshefte f{\"u}r Mathematik und Physik}, 38:\penalty0
  173--198, 1931.
\newblock \doi{10.1007/BF01700692}.

\bibitem[Green et~al.(2007)Green, Karvounarakis, and
  Tannen]{GreenKarvounarakisTannen2007}
Todd~J. Green, Grigoris Karvounarakis, and Val Tannen.
\newblock Provenance semirings.
\newblock In \emph{Proceedings of the 26th ACM SIGMOD-SIGACT-SIGART Symposium
  on Principles of Database Systems (PODS '07)}, pages 31--40. ACM, 2007.
\newblock \doi{10.1145/1265530.1265535}.

\bibitem[Grothendieck(1971)]{GrothendieckSGA1}
Alexander Grothendieck.
\newblock \emph{Rev{\^e}tements {\'e}tales et groupe fondamental ({SGA} 1)},
  volume 224 of \emph{Lecture Notes in Mathematics}.
\newblock Springer, 1971.
\newblock \doi{10.1007/BFb0058656}.

\bibitem[Hoare(1969)]{Hoare1969}
C.~A.~R. Hoare.
\newblock An axiomatic basis for computer programming.
\newblock \emph{Communications of the ACM}, 12\penalty0 (10):\penalty0
  576--580, 1969.
\newblock \doi{10.1145/363235.363259}.

\bibitem[Janelidze and Tholen(1994)]{JanelidzeTholen1994}
G.~Janelidze and W.~Tholen.
\newblock Facets of descent, {I}.
\newblock \emph{Applied Categorical Structures}, 2\penalty0 (3):\penalty0
  245--281, 1994.
\newblock \doi{10.1007/BF00878100}.

\bibitem[Kemeny and Snell(1960)]{KemenySnell1960}
John~G. Kemeny and J.~Laurie Snell.
\newblock \emph{Finite {M}arkov Chains}.
\newblock University Series in Undergraduate Mathematics. Van Nostrand, 1960.
\newblock ISBN 0442043287.
\newblock Reprinted by Springer, 1976.

\bibitem[Ladyman and Ross(2007)]{LadymanRoss2007}
James Ladyman and Don Ross.
\newblock \emph{Every Thing Must Go: Metaphysics Naturalized}.
\newblock Oxford University Press, 2007.
\newblock \doi{10.1093/acprof:oso/9780199276196.001.0001}.

\bibitem[Libkin(2004)]{Libkin2004}
Leonid Libkin.
\newblock \emph{Elements of Finite Model Theory}.
\newblock Springer, 2004.
\newblock \doi{10.1007/978-3-662-07003-1}.

\bibitem[L{\"o}b(1955)]{Lob1955}
M.~H. L{\"o}b.
\newblock Solution of a problem of {L}eon {H}enkin.
\newblock \emph{The Journal of Symbolic Logic}, 20\penalty0 (2):\penalty0
  115--118, 1955.
\newblock \doi{10.2307/2266895}.

\bibitem[Mac~Lane(1998)]{MacLane1998}
Saunders Mac~Lane.
\newblock \emph{Categories for the Working Mathematician}, volume~5 of
  \emph{Graduate Texts in Mathematics}.
\newblock Springer, 2 edition, 1998.
\newblock \doi{10.1007/978-1-4757-4721-8}.

\bibitem[Martin-L{\"o}f(1984)]{MartinLof1984}
Per Martin-L{\"o}f.
\newblock \emph{Intuitionistic Type Theory}, volume~1 of \emph{Studies in Proof
  Theory. Lecture Notes}.
\newblock Bibliopolis, 1984.
\newblock ISBN 9788870881059.
\newblock Notes by Giovanni Sambin of lectures given in Padua, June 1980.

\bibitem[Merkle(1988)]{Merkle1987}
Ralph~C. Merkle.
\newblock A digital signature based on a conventional encryption function.
\newblock In Carl Pomerance, editor, \emph{Advances in Cryptology --- CRYPTO
  '87}, volume 293 of \emph{Lecture Notes in Computer Science}, pages 369--378.
  Springer, 1988.
\newblock \doi{10.1007/3-540-48184-2_32}.

\bibitem[Milner(1989)]{Milner1989}
Robin Milner.
\newblock \emph{Communication and Concurrency}.
\newblock Prentice Hall International Series in Computer Science. Prentice
  Hall, 1989.
\newblock ISBN 9780131150072.

\bibitem[Nakahara(2003)]{Nakahara2003}
Mikio Nakahara.
\newblock \emph{Geometry, Topology and Physics}.
\newblock Graduate Student Series in Physics. Institute of Physics Publishing,
  2 edition, 2003.
\newblock ISBN 9780750306065.

\bibitem[Pearl(2009)]{Pearl2009}
Judea Pearl.
\newblock \emph{Causality: Models, Reasoning, and Inference}.
\newblock Cambridge University Press, 2 edition, 2009.
\newblock \doi{10.1017/CBO9780511803161}.

\bibitem[Seifert(2012)]{Seifert2012}
Udo Seifert.
\newblock Stochastic thermodynamics, fluctuation theorems and molecular
  machines.
\newblock \emph{Reports on Progress in Physics}, 75\penalty0 (12):\penalty0
  126001, 2012.
\newblock \doi{10.1088/0034-4885/75/12/126001}.

\bibitem[Tsiokos(2026{\natexlab{a}})]{Tsiokos2026SixBirdsFoundations}
Ioannis Tsiokos.
\newblock Six birds: Foundations of emergence calculus, 2026{\natexlab{a}}.
\newblock arXiv:2602.00134.

\bibitem[Tsiokos(2026{\natexlab{b}})]{Tsiokos2026SixBirdsNoGo}
Ioannis Tsiokos.
\newblock Six birds: No-go theorems for audited emergence, 2026{\natexlab{b}}.
\newblock Zenodo preprint.

\bibitem[Tsiokos(2026{\natexlab{c}})]{Tsiokos2026SixBirdsProtocolTrap}
Ioannis Tsiokos.
\newblock Six birds protocol trap: Holonomy without entropy production,
  2026{\natexlab{c}}.
\newblock Zenodo preprint.

\bibitem[{van Fraassen}(1980)]{vanFraassen1980}
Bas~C. {van Fraassen}.
\newblock \emph{The Scientific Image}.
\newblock Oxford University Press, 1980.
\newblock \doi{10.1093/0198244274.001.0001}.

\end{thebibliography}

\end{document}
